\documentclass{template}

\usepackage{mathtools}
\usepackage{mdframed}
\usepackage{varwidth} 
\pgfplotsset{compat=1.18}

\newcommand{\qedwhite}{\hfill \ensuremath{\Box}}

% Random sampling
\newcommand\leftarrowS{\leftarrow\joinrel\smalldollar}
\newcommand\rightarrowS{\smalldollar\joinrel\rightarrow}

\makeatletter
\newcommand{\smalldollar}{\mathrel{\mathpalette\small@dollar\relax}}
\newcommand{\small@dollar}[2]{%
  \vcenter{\hbox{%
    $#1\textnormal{\fontsize{0.7\dimexpr\f@size pt}{0}\selectfont\$}$%
  }}%
}
\makeatother


\begin{document}

\begin{center}
    \Large\textcolor{blue!70!black}{\textbf{Introduction to Lattice-based Cryptography}} \\
    \normalsize\textit{Dmytro Zakharov}
\end{center}

\begin{abstract}
    In recent years, the interest in lattice-based cryptography has been
    steadily increasing. One primary reason is, of course, that problems related
    to lattices and their derivatives are widely believed to be \textit{hard
    even for quantum computers}. Furthermore, there are settings where the
    algebraic structure of lattices is natural: e.g., for fully-homomorphic
    encryption.

    This brief paper is a technical introduction to lattice-based cryptography.
    We give formal definitions of lattices and introduce key related hardness
    assumptions such as SVP, CVP, and several variants of LWE. We show how to
    build public-key encryption (PKE) and digital signature schemes based on LWE
    assumptions and partially prove the security of key constructions. Finally,
    we provide a high-level specification of Kyber PKE and Dilithium signature
    schemes.
    
    We aim to make these notes accessible while preserving rigorousness where
    appropriate. Basic knowledge of cryptography and algebra is required.
\end{abstract}

\tableofcontents

\newpage

\section{Lattices}\label{section:lattices}

\subsection*{Notation and Basic Notions}

All vectors are written in bold (e.g., $\mathbf{x}$ or $\mathbf{y}$) and
considered \textit{column} vectors by default. All real vector spaces are
Euclidean: that is, for every $\mathbf{x}, \mathbf{y} \in V \subseteq
\mathbb{R}^d$, we define the inner product $\langle \cdot, \cdot \rangle$ and
$\ell_2$ norm $\|\cdot\|$ in a standard fashion:
\begin{equation*}
    \langle \mathbf{x}, \mathbf{y} \rangle \triangleq \sum_{i \in [d]} x_iy_i, \quad \|\mathbf{x}\| \triangleq \sqrt{\langle \mathbf{x}, \mathbf{x} \rangle},
\end{equation*}

where we use $[d]$ to denote the set $\{1,\dots,d\}$. Additionally, we use $(d)$
to denote the ranging set using 0-indexed notation: $\{0,\dots,d-1\}$. Other
norms will be introduced on the go when needed. Finally, we use notation $x
\leftarrowS \mathcal{X}$ to denote sampling from the distribution $\mathcal{X}$.
If $\mathcal{X}$ is a set, the underlying distribution is uniform over
$\mathcal{X}$. 

\subsection{Definition}

We start our discussion with the definition of the lattice. Following
\cite{nist}, we provide two definitions: one is rather algebraic, while the
other is more intuitive and practical. In fact, both definitions will be helpful
since some facts are sometimes easier to prove using the first definition. 
\begin{definition}\label{def:lattice}
    Let $V$ be a real Euclidean vector space of finite dimension $d$. The
    following two definitions of \textbf{lattice} $\Lambda \subseteq V$ of rank
    $r \leq d$ are equivalent:
    \begin{enumerate}
        \item[\textit{(i)}] $\Lambda$ is a discrete additive subgroup of $V$
        (with $\Lambda \cong \mathbb{Z}^r$ as $\mathbb{Z}$-modules).
        \item[\textit{(ii)}] $\Lambda = \mathcal{L}(B) \triangleq \{B\mathbf{z}:
        \mathbf{z} \in \mathbb{Z}^d\}$ for some matrix $B \in \mathbb{R}^{r
        \times d}$ where its columns, further denoted as
        $\{\mathbf{b}_1,\dots,\mathbf{b}_r\} \subseteq \mathbb{R}^d$, are
        linearly independent over $V$. 
    \end{enumerate}
\end{definition}

While the second definition is clear, the first one requires some deciphering.
Since we work in the Euclidean space, it makes sense to define what it means for
the set $S \subseteq V$ to be \textit{discrete}. Intuitively, it should mean
that any two points $\mathbf{s}_1,\mathbf{s}_2 \in S$ are separated by some
``considerable'' distance. To translate it into the mathematical language, we
require that:
\begin{equation*}
    \inf_{\mathbf{s}_1 \neq \mathbf{s}_2 \in S} \|\mathbf{s}_1-\mathbf{s}_2\| > 0.
\end{equation*}

Besides $\Lambda$ being discrete, additive subgroup structure of $\Lambda$
guarantees that $\Lambda$ is closed under any $\mathbb{Z}$-linear combinations:
that is, for any $\mathbf{x}, \mathbf{y} \in \Lambda$ and $\lambda, \mu \in
\mathbb{Z}$, we have $\lambda\mathbf{x} + \mu\mathbf{y} \in \Lambda$. Showing
this property might sometimes be easier than finding the matrix $B$ in
definition~\textit{(ii)} explicitly.

\textbf{Geometric Intuition.} While definition~\textit{(i)} is useful, the
definition~\textit{(ii)} gives an intuitive way of thinking of lattice $\Lambda
= \mathcal{L}(B)$. If vector space $V$ is typically thought as a set of all
$\mathbb{R}$-linear combinations\footnote{We further use notation $S+S':=\{s+s':
s \in S, s' \in S'\}$ and for $a \in S$, set $aS := \{as: s \in S\}$.} $\sum_{i
\in [d]}\mathbf{e}_i\mathbb{R}$ over some basis $\{\mathbf{e}_i\}_{i \in [d]}$,
the lattice $\Lambda$ can be thought as all $\mathbb{Z}$-linear combinations
$\sum_{i \in [r]}\mathbf{b}_i\mathbb{Z}$ over real basis $\{\mathbf{b}_i\}_{i
\in [r]} \subseteq \mathbb{R}^d$. See \Cref{fig:lattice} for illustration.

\begin{remark}\label{remark:lattice-in-integers}
Note that the basis of lattice $\Lambda$ is not required to lie in
$\mathbb{Z}^d$, compared to coefficients of the linear combination. However, in
practice they usually do: note that we still can represent basis in memory only
in the rational form (that is, $\mathbf{b}_i \in \mathbb{Q}^d$). With the proper
rescaling, representation in $\mathbb{Z}^d$ would suffice.
\end{remark}

\begin{remark}
    In this section we consider \textit{full-rank} lattices only: that is,
    $\Lambda \cong \mathbb{Z}^{d}$ for $d=\dim(V)$.
\end{remark}

\begin{figure}[ht]
  \centering
  \begin{tikzpicture}
    \coordinate (Origin)   at (0,0);
    \coordinate (XAxisMin) at (-2,0);
    \coordinate (XAxisMax) at (5,0);
    \coordinate (YAxisMin) at (0,-2);
    \coordinate (YAxisMax) at (0,5);
    \draw [ultra thick, gray,-stealth] (XAxisMin) -- (XAxisMax);
    \draw [ultra thick, gray,-stealth] (YAxisMin) -- (YAxisMax);

    \clip (-2,-2) rectangle (5cm,5cm); 
    \pgftransformcm{1}{0.6}{0.7}{1}{\pgfpoint{0cm}{0cm}}

    \coordinate (Bone) at (0,2);
    \coordinate (Btwo) at (2,-2);
    \draw[style=help lines, dashed, ultra thick, black!20] (-14,-14) grid[step=2cm] (14,14);

    \foreach \x in {-7,-6,...,7}{
      \foreach \y in {-7,-6,...,7}{
        \node[draw,circle,inner sep=2pt,fill] at (2*\x,2*\y) {};
      }
    }

    \draw [ultra thick,-latex,blue!70!black] (Origin)
        -- (Bone) node [right] {$\mathbf{b}_1$};
    \draw [ultra thick,-latex,blue!70!black] (Origin)
        -- (Btwo) node [below right] {$\mathbf{b}_2$};
    \filldraw[fill=red!70!purple, dashed, very thick, fill opacity=0.3, draw=black] 
        (Origin) -- (Bone) -- ($(Bone) + (Bone) - (Btwo)$) -- ($(Bone)-(Btwo)$) -- cycle;
  
    %\draw [ultra thick,-latex,red!70!purple] (Origin)
    %    -- ($0+1*(Bone)-3*(Btwo)$) node [above left] {$\mathbf{b}_1-3\mathbf{b}_2$};
    \draw [ultra thick,-latex,red!70!purple] (Origin)
        -- ($(Bone)-(Btwo)$) node [above left] {$\mathbf{b}_1-\mathbf{b}_2$};
    \filldraw[fill=blue!70!black, dashed, very thick, fill opacity=0.3, draw=black] 
    (Origin) -- (Bone) -- ($(Bone)+(Btwo)$) -- (Btwo) -- cycle;
    \node at ($0.5*(Bone)+0.5*(Btwo)$) {\footnotesize$\mathcal{P}(\textcolor{blue!70!black}{B})$};
    \node at ($0.5*(Bone) + 0.5*($(Bone)-(Btwo)$)$) {\footnotesize$\;\mathcal{P}(\textcolor{red!70!purple}{B'})$};
  \end{tikzpicture}
  \caption{Lattice $\Lambda=\mathcal{L}(\textcolor{blue!70!black}{B})$ of rank
  $2$ with basis $\textcolor{blue!70!black}{B} =
  \{\textcolor{blue!70!black}{\mathbf{b}_1},\textcolor{blue!70!black}{\mathbf{b}_2}\}$.
  $\textcolor{blue!70!black}{B}$ is an example of a ``good'' basis where basis
  vectors are almost orthogonal. The basis $\textcolor{red!70!purple}{B'} =
  \{\textcolor{blue!70!black}{\mathbf{b}_1},\textcolor{red!70!purple}{\mathbf{b}_1-\mathbf{b}_2}\}$
  of the same lattice $\Lambda=\mathcal{L}(\textcolor{red!70!purple}{B'})$ is ``worse'' (the vectors seem more parallel and
  fundamental domain $\mathcal{P}(\textcolor{red!70!purple}{B'})$ is skewed).}
  \label{fig:lattice}
\end{figure}

\subsection{Fundamental Measures}

\subsubsection{Fundamental Domain}

A basis $B$ of $\Lambda = \mathcal{L}(B)$ is associated with the fundamental
domain $\mathcal{P}(B) \triangleq \sum_{i \in
[d]}\mathbf{b}_i[0,1)$\footnote{Sometimes, the interval $[-1/2, 1/2)$ is used
instead of $[0,1)$ to center $\mathcal{P}(B)$ with respect to origin: in such
case, $\mathcal{P}(B)$ is called a \textit{centered parallelepiped of $B$}}.

What is special about this parallelepiped is that it represents all elements in
the quotient group $V/\Lambda$: any element $\mathbf{x} \in V$ can be uniquely
written as $\mathbf{x} = \mathbf{v} + \mathbf{p}$ where $\mathbf{v} \in \Lambda$
and $\mathbf{p} \in \mathcal{P}(B)$. Moreover, the shape of this
$\mathcal{P}(B)$ determines the ``quality'' of the basis $B$: if
$\mathcal{P}(B)$ is skewed and is not concentrated in the origin, it is
typically considered ``bad'' as some computational problems become much harder
to solve. In turn, a non-skewed basis $B'$ primarily concentrated around the
origin is considered ``good'' as $\{\mathbf{b}_i'\}_{i \in [d]}$ much more
closely resembles the orthogonal basis than one of $B$ (and of course, the
orthogonal basis in a sense is ``perfect'' because it is very ``nice'' to work
with). Again, see \Cref{fig:lattice} for illustration.

Another natural quantity to consider is the volume of such parallelepiped,
called \textit{covolume}.

\begin{definition}[Covolume]
    The \textbf{covolume} $\text{Vol}(\Lambda)$ of lattice $\Lambda$ is defined as:
    \begin{equation*}
        \text{Vol}(\Lambda) \triangleq \text{Vol}(V/\Lambda) = \sqrt{|\det(B^{\top}B)|}.
    \end{equation*}
\end{definition}

Note that the definition of $\mathcal{P}(B)$ depends on the choice of $B$ since
the basis of $\Lambda$ is generally not unique. Thus, it might seem that
covolume should also depend on the choice of $B$. However, we show the
following.

\begin{proposition}
    The \textbf{covolume} $\text{Vol}(\Lambda)$ is an invariant of $\Lambda$:
    that is, it does not depend on the choice of $B$.
\end{proposition}

\textbf{Proof.} Suppose $\Lambda = \mathcal{L}(B_0) = \mathcal{L}(B_1)$. Since
$B_0$ and $B_1$ both generate $\Lambda$, one has $B_1 = UB_0$ for some matrix $U
\in \text{GL}_d(\mathbb{Z})$. Since $U$ is an integer matrix with integer
inverse, it is unimodular: that is, $|\det U| = 1$. It is then easy to see
$|\det B_1^{\top}B_1| = |\det B_0^{\top}U^{\top}UB_0| = |\det B_0^{\top}B_0| =
\text{Vol}(\Lambda)^2$. \qedwhite

\subsubsection{Successive Minima}

Turns out that one of the hardest computational problems is computing the
shortest vectors of $\Lambda$. However, we need some preliminary work to
formalize this statement. 

Denote $B_{\rho}(\mathbf{x}_0) := \{\mathbf{x} \in V:
\|\mathbf{x}-\mathbf{x}_0\|<\rho\}$ --- an open ball in $V$ of radius $\rho$
centered at $\mathbf{x}_0$. 

\begin{definition}[Successive minima]
    The \textbf{$i^{\text{th}}$ minimum $\lambda_i(\Lambda)$} is the radius 
    of the smallest sphere centered at the origin containing $i$ linearly 
    independent lattice vectors:
    \begin{equation*}
        \lambda_i(\Lambda) \triangleq \inf \{\rho \in \mathbb{R}_{\geq 0}: \dim(\text{span}(\Lambda \cap B_{\rho}(\mathbf{0}))) \geq i\}.
    \end{equation*}
\end{definition}

In particular, $\lambda_1(\Lambda) = \min_{\mathbf{x} \neq \mathbf{y} \in
\Lambda}\|\mathbf{x}-\mathbf{y}\| = \min_{\mathbf{x} \in
\Lambda\setminus\{\mathbf{0}\}}\|\mathbf{x}\|$ is called the \textit{first
minimum} while the vector of length $\lambda_1(\Lambda)$ --- the shortest
vector. The natural question to ask though is whether it exists at all: after
all, in $\mathbb{R}^d$, for instance, it is clearly undefined. We give
affirmative answer.

\subsubsection*{Gram-Schmidt Orthogonalization*}

This is probably a good place to remind the reader of the Gram-Schmidt
Orthogonalization. Recall that the main goal of Gram-Schmidt Orthogonalization
is, based on sequence of vectors $\{\mathbf{b}_1,\dots,\mathbf{b}_n\}$, produce
orthogonal vectors $\{\mathbf{b}_1^*,\dots,\mathbf{b}_n^*\}$ such that
$\text{span}\{\mathbf{b}_i\}_{i \in [n]} = \text{span}\{\mathbf{b}_i^*\}_{i \in
[n]}$.

For that, set $\mathbf{b}_1^* := \mathbf{b}_1$. Now we shall produce the vector
$\mathbf{b}_2^*$ that is orthogonal to $\mathbf{b}_1$. For that, we set
$\mathbf{b}_2^*$ to the component of $\mathbf{b}_2$ that is orthogonal to
$\mathbf{b}_1$. Recall that the angle $\theta$ between $\mathbf{b}_1$ and
$\mathbf{b}_2$ can be found as $\cos\theta = \langle \mathbf{b}_1, \mathbf{b}_2
\rangle /(\|\mathbf{b}_1\| \cdot \|\mathbf{b}_2\|)$. The component of
$\mathbf{b}_2$ parallel to $\mathbf{b}_1$ can be thus found as
$\widehat{\mathbf{b}}_1 \cdot \|\mathbf{b}_2\|\cos\theta = \langle \mathbf{b}_1,
\mathbf{b}_2 \rangle\mathbf{b}_1 /\|\mathbf{b}_1\|^2$ and so $\mathbf{b}_2^* :=
\mathbf{b}_1 - \langle \mathbf{b}_1, \mathbf{b}_2 \rangle\mathbf{b}_1
/\|\mathbf{b}_1\|^2$.

In a similar fashion, one obtains the general relations for the rest of vectors:
\begin{equation*}
    \mathbf{b}_i^* := \mathbf{b}_i - \sum_{j=0}^{i-1}\mu_{i,j}\mathbf{b}_j^*, \quad \mu_{i,j} := \frac{\langle \mathbf{b}_i, \mathbf{b}_j^*\rangle}{\langle \mathbf{b}_j^*, \mathbf{b}_j^*\rangle}.
\end{equation*}

We shall further denote the result of Gram-Schmidt orthogonalization of $B$ by
$B^*$. The question that might naturally arise: given such orthogonalization,
can't we then always turn a ``bad'' basis of lattice $\Lambda$ into a ``good''
basis? Unfortunately, typically $\mathcal{L}(B) \neq \mathcal{L}(B^*)$: the
orthogonalization of a lattice basis is (typically) not a lattice basis.
However, the idea of Gram-Schmidt orthogonalization lies in the core of many
attacks on lattices. Besides, it has several nice properties which we list
below:
\begin{proposition}\label{prop:gram-schmidt-props}
    Suppose $\Lambda = \mathcal{L}(B)$ is a lattice and $B^*$ is the
    Gram-Schmidt orthogonalization of $B$ with columns $\{\mathbf{b}^*_j\}_{j
    \in [d]}$. Then:
    \begin{enumerate}
        \item[\textit{(a)}] $\text{Vol}(\Lambda) = \prod_{j=1}^d \|\mathbf{b}_j^*\|$.
        \item[\textit{(b)}] $\lambda_1(\Lambda) \geq \min_{j \in [d]}\|\mathbf{b}_j^*\| >
        0$. 
    \end{enumerate}
\end{proposition}

\textbf{Proof.} 

\textit{Part (a).} Note that $\det(B) = \det(B^*)$: the process of Gram-Schmidt
orthogonalization only comprises adding one row to another with some factor
which does not change the determinant. Then, $\text{Vol}(\Lambda) =
\sqrt{\det(B^{\top}B)} = \sqrt{\det((B^*)^{\top}B^*)} = \sqrt{\det
\text{diag}\{\|\mathbf{b}_1^*\|^2, \dots, \|\mathbf{b}_d^*\|^2\}} = \prod_{j \in
[d]}\|\mathbf{b}_j^*\|$.

\textit{Part (b)} (adapted from~\cite{complexity-goldwasser}). Consider a
generic non-zero lattice vector $B\mathbf{x}$ ($\mathbf{x} \in \mathbb{Z}^d
\setminus \{\mathbf{0}\}$) and let $i$ be the biggest index such that $x_i \neq
0$. We first prove that $|\langle B\mathbf{x}, \mathbf{b}_i^* \rangle| \geq
\|\mathbf{b}_i^*\|^2$. From the definition of $i$, we know $B\mathbf{x} =
\sum_{j=1}^i \mathbf{b}_jx_j$. Thus:
\begin{align*}
    \langle B\mathbf{x}, \mathbf{b}_i^* \rangle = \sum_{j=1}^i \langle \mathbf{b}_j, \mathbf{b}_i^*\rangle x_j &= \langle \mathbf{b}_i,\mathbf{b}_i^* \rangle x_i\\
    &= \langle \mathbf{b}_i^* + \sum_{j<i}\mu_{i,j}\mathbf{b}_j^*, \mathbf{b}_i^*\rangle x_i \\
    &= \langle \mathbf{b}_i^*, \mathbf{b}_i^*\rangle x_i + \sum_{j < i}\mu_{i,j}\langle \mathbf{b}_j^*, \mathbf{b}_i^*\rangle x_i \\
    &= \|\mathbf{b}_i^*\|^2x_i.
\end{align*}
From this follows $|\langle B\mathbf{x}, \mathbf{b}_i^* \rangle| \geq
\|\mathbf{b}_i^*\|^2|x_i| \geq \|\mathbf{b}_i^*\|^2$ and therefore
$\|B\mathbf{x}\|\cdot\|\mathbf{b}_i^*\| \geq \|\mathbf{b}_i^*\|^2$. Clearly
$\|\mathbf{b}_i^*\| \neq 0$ since $\mathbf{b}_i$'s are linearly independent, so
$\|B\mathbf{x}\| \geq \|\mathbf{b}_i^*\| \geq \min_{j \in
[d]}\|\mathbf{b}_j^*\|$. By definition of $\lambda_1(\Lambda)$, the result
follows. \qedwhite

\begin{proposition}
    For lattice $\Lambda$, there exists $\mathbf{v} \in \Lambda$ such that
    $\|\mathbf{v}\| = \lambda_1(\Lambda)$.
\end{proposition}

\textbf{Proof.} By the definition of $\lambda_1(\Lambda)$, there exists a
sequence $\{\mathbf{v}_n\}_{n \in \mathbb{N}} \subseteq \Lambda$ such that
$\lim_{n \to \infty}\|\mathbf{v}_n\| = \lambda_1(\Lambda)$. From
Proposition~\ref{prop:gram-schmidt-props}, $\lambda_1(\Lambda) > 0$, so for all
sufficiently large $n$ it must be $\|\mathbf{v}_n\| \leq 2\lambda_1(\Lambda)$,
i.e., lattice vector $\mathbf{v}_n$ belongs to the closed ball
$\overline{B}_{2\lambda_1(\Lambda)}(\mathbf{0})$. But this set is compact, so we
might extract a convergent subsequence $\{\mathbf{v}_{i_j}\}_{j \in \mathbb{N}}$
with limit $\mathbf{w} = \lim_{j \to \infty}\mathbf{v}_{i_j}$ and
$\|\mathbf{w}\| = \lambda_1(\Lambda)$. We claim that $\mathbf{w} \in \Lambda$.
By definition of $\mathbf{w}$, $\lim_{j \to \infty}\|\mathbf{v}_{i_j} -
\mathbf{w}\|=0$. Therefore for all sufficiently large $j$, $\|\mathbf{v}_{i_j} -
\mathbf{w}\| < \lambda_1(\Lambda)/2$. By triangle inequality for sufficiently
large $j$ and all $k>j$, $\|\mathbf{v}_{i_j}-\mathbf{v}_{i_k}\| \leq
\|\mathbf{v}_{i_j} - \mathbf{w}\| + \|\mathbf{w} - \mathbf{v}_{i_k}\| <
\lambda_1(\Lambda)$. But $\mathbf{v}_{i_j}-\mathbf{v}_{i_k} \in \Lambda$, and 
no non-zero lattice vector can have length strictly less than $\lambda_1(\Lambda)$.
Thus, $\mathbf{v}_{i_k}=\mathbf{v}_{i_j}$ for all $k>j$. Thus $\mathbf{w}=\mathbf{v}_{i_j} \in \Lambda$. \qedwhite

\subsection{Minkowski's Theorem}

Clearly, one expects the dependence between the covolume of a lattice
$\text{Vol}(\Lambda)$ and the first minimum $\lambda_1(\Lambda)$. In particular,
for smaller covolume we expect the smaller first minimum $\lambda_1(\Lambda)$
(or alternatively, one can say the lattice with large covolume cannot have a
small first minimum). Moreover, by dimensionality analysis one should expect
that $\lambda_1(\Lambda) \propto \text{Vol}(\Lambda)^{1/d}$. The Minkowski
theorem gives the inequality between these two quantities.

\begin{theorem}[Minkowski's Theorems]\label{theorem:minkowski}
    Let $\Lambda=\mathcal{L}(B)$ be the lattice of rank $d$ with covolume
    $\text{Vol}(\Lambda) = \det B$. Then,
    \begin{enumerate}[align=left]
        \item[\textbf{First Minkowski's Theorem}] $\lambda_1(\Lambda) < \sqrt{d}\text{Vol}(\Lambda)^{1/d}$.
        \item[\textbf{Second Minkowski's Theorem}] $\prod_{i \in [d]}\lambda_i^{1/d} < \sqrt{d}\text{Vol}(\Lambda)^{1/d}$. 
    \end{enumerate}
\end{theorem}

\textbf{Proof.} See \cite[Section 1.3]{complexity-goldwasser}.

While the Minkowski's Theorem turns out to be asymptotically tight for ``bad''
lattices, in general $\lambda_1(\Lambda)$ might be much smaller than
$\sqrt{d}\text{Vol}(\Lambda)^{1/d}$. Consider a very simple example: let $d=2$
and $\mathbf{b}_1 := \varepsilon \mathbf{e}_1$, $\mathbf{b}_2 :=
(1/\varepsilon)\mathbf{e}_2$ (where $\mathbf{e}_1,\mathbf{e}_2$ is the canonical
basis for $\mathbb{R}^2$). Minkowski's Theorem gives $\lambda_1(\Lambda) \leq
\sqrt{2}$, however one easily sees that $\lambda_1(\Lambda) = \varepsilon$ can
be made arbitrarily small.

Moreover, Minkowski's Theorem proof is not constructive: we know that a vector
of length below $\sqrt{d}\text{Vol}(\Lambda)^{1/d}$ exists, but there is no
concrete algorithm to find it, let alone $\lambda_1(\Lambda)$. This became the
core of one of the most well-known lattice computational problem: finding
shortest vector in $\Lambda$. We introduce this and related problems in the next
subsection.

\subsection{Lattice Computational Problems}

In this section, we finally define some problems that are currently believed to
be unsolvable in polynomial time (under reasonable choice of parameters, of
course). 

\subsubsection{Shortest Vector Problem (SVP)}

As previously announced, we start with the \textit{shortest vector
problem} (SVP), which requires to find the vector of length
$\lambda_1(\Lambda)$.

\begin{definition}[Lattice Generation Algorithm]
    A \textbf{lattice generation algorithm} $\mathsf{LGen}(1^{\mu})$ is a
    probabilistic polynomial time (PPT) algorithm that takes the security
    parameter $1^{\mu}$ and outputs the lattice basis\footnote{As previously
    noted in Remark~\ref{remark:lattice-in-integers}, we might assume that $B$
    is given in the integer representation.} $B \in \mathbb{Z}^{\mu \times \mu}$
    where the lattice $\Lambda=\mathcal{L}(B)$ is of rank $\mu$.
\end{definition}

\begin{definition}[Shortest Vector Problem]\label{def:svp}
    For a given adversary $\mathcal{A}$ and lattice generation algorithm
    $\mathsf{LGen}(1^{\mu})$, define the game
    $\underline{\mathsf{Game}\,\mathsf{SVP}}^{\mathcal{A}}_{\mathsf{LGen}}$ as follows:
    \begin{enumerate}
        \item Sample the lattice basis $B \gets \mathsf{LGen}(1^{\mu})$ and send to $\mathcal{A}$;
        \item $\mathcal{A}$ outputs some $\mathbf{x} \in \mathbb{Z}^{\mu} \setminus \{\mathbf{0}\}$ based on $B \in \mathbb{Z}^{\mu \times \mu}$;
        \item \textbf{return} $1$ iff $\|B\mathbf{x}\|=\lambda_1(\Lambda)$ and $0$ otherwise.
    \end{enumerate}
    We define the $\mathcal{A}$'s \textbf{advantage in solving the shortest
    vector problem} as:
    \begin{equation*}
        \mathsf{Adv}^{\text{SVP}}_{\mathcal{A},\mathsf{LGen}}(\mu) := \text{Pr}\left[\underline{\mathsf{Game}\,\mathsf{SVP}}^{\mathcal{A}}_{\mathsf{LGen}}(\mu) = 1\right]
    \end{equation*} 
\end{definition}

\begin{definition}[SVP Assumption]
    We say that \textbf{Shortest Vector Problem (SVP) assumption} holds if for 
    any PPT adversary $\mathcal{A}$,
    \begin{equation*}
        \mathsf{Adv}^{\text{SVP}}_{\mathcal{A},\mathsf{LGen}}(\mu) = \text{negl}(\mu).
    \end{equation*}
\end{definition}

\subsubsection{Approximate Shortest Vector Problem ($\text{SVP}_{\gamma}$)}

Turns out that the setting of \Cref{def:svp} is too strict: finding the shortest
non-zero vector in lattice $\Lambda$ is too hard. Moreover, one expects that
building cryptographic protocols on this assumption solely would be highly hard.
So instead, we will require not finding the shortest vector, but
``short-enough'', where the natural measure of ``short'' is ``the multiple of
$\lambda_1(\Lambda)$''.

\begin{definition}[Approximate Shortest Vector Problem]\label{def:appr-svp}
    For a given adversary $\mathcal{A}$ and lattice generation algorithm
    $\mathsf{LGen}(1^{\mu})$, define the game
    $\underline{\mathsf{Game}\,\mathsf{SVP}}^{\mathcal{A}}_{\textcolor{red!70!purple}{\gamma},\mathsf{LGen}}$ as follows:
    \begin{enumerate}
        \item Sample the lattice basis $B \gets \mathsf{LGen}(1^{\mu})$ and send to $\mathcal{A}$;
        \item $\mathcal{A}$ outputs some $\mathbf{x} \in \mathbb{Z}^{\mu} \setminus \{\mathbf{0}\}$ based on $(B,\textcolor{red!70!purple}{\gamma})$;
        \item \textbf{return} $1$ iff $\|B\mathbf{x}\| \leq \textcolor{red!70!purple}{\gamma}\lambda_1(\Lambda)$ and $0$ otherwise.
    \end{enumerate}
    We define the $\mathcal{A}$'s \textbf{advantage in solving the approximate
    shortest vector problem} as:
    \begin{equation*}
        \mathsf{Adv}^{\text{SVP}_{\textcolor{red!70!purple}{\gamma}}}_{\mathcal{A},\mathsf{LGen}}(\mu) := \text{Pr}\left[\underline{\mathsf{Game}\,\mathsf{SVP}}^{\mathcal{A}}_{\textcolor{red!70!purple}{\gamma},\mathsf{LGen}}(\mu) = 1\right]
    \end{equation*} 
\end{definition}

\begin{definition}[Approximate SVP Assumption]\label{def:approximate-svp-assumption}
    We say that \textbf{Approximate Shortest Vector Problem (SVP) assumption}
    with given $\gamma$ holds if for any PPT adversary $\mathcal{A}$,
    \begin{equation*}
        \mathsf{Adv}^{\text{SVP}_{\gamma}}_{\mathcal{A},\mathsf{LGen}}(\mu) = \text{negl}(\mu).
    \end{equation*}
\end{definition}

Typically, $\gamma$ is analyzed as a function of the rank of the lattice $\mu$.
According to \cite[Section 2.2.5]{nist}, for example, vectors of length below
$\gamma\lambda_1(\Lambda)$ for $\gamma = \sqrt{\mu}$ are considered ``very short''
while $\gamma = \mathsf{poly}(\mu)$ is the most interesting case. As one
expects, the Approximate SVP assumption does not hold for $\gamma=2^{\mu}$ by
means of Lenstra-Lenstra-Lovasz (LLL) algorithm.

\subsubsection{Approximate Closest Vector Problem ($\text{CVP}_{\gamma}$)}\label{section:approximate-cvp}

A somewhat similar problem to SVP is, given some target vector $\mathbf{t} \in
\mathbb{Z}^{\mu}$, find the closest vector to $\mathbf{t}$ on the lattice. This
problem also turns out to be computationally hard. Similarly to SVP, the raw CVP
problem is too strict, so we give formal definition for approximate CVP straight away.

\begin{definition}[Approximate Closest Vector Problem]\label{def:cvp}
    For a given adversary $\mathcal{A}$ and lattice generation algorithm
    $\mathsf{LGen}(1^{\mu})$, define the game
    $\underline{\mathsf{Game}\,\mathsf{CVP}}^{\mathcal{A}}_{\gamma,\mathsf{LGen}}$ as follows:
    \begin{enumerate}
        \item Sample the lattice basis $B \gets \mathsf{LGen}(1^{\mu})$;
        \item Sample target vector $\mathbf{t} \leftarrowS \mathbb{Z}^{\mu}$ and
        send $(B,\mathbf{t})$ to $\mathcal{A}$;
        \item $\mathcal{A}$ outputs some $\mathbf{x} \in \mathbb{Z}^{\mu}$ based
        on $(B,\mathbf{t},\gamma)$;
        \item \textbf{return} $1$ iff $\|B\mathbf{x}-\mathbf{t}\| \leq \gamma\|B\mathbf{y}-\mathbf{t}\|$
        for all $\mathbf{y} \neq \mathbf{x}$ and $0$ otherwise.
    \end{enumerate}
    We define the $\mathcal{A}$'s \textbf{advantage in solving the approximate closest
    vector problem} as:
    \begin{equation*}
        \mathsf{Adv}^{\text{CVP}_{\gamma}}_{\mathcal{A},\mathsf{LGen}}(\mu) := \text{Pr}\left[\underline{\mathsf{Game}\,\mathsf{CVP}}^{\mathcal{A}}_{\gamma,\mathsf{LGen}}(\mu) = 1\right]
    \end{equation*} 
\end{definition}

Similarly to~\Cref{def:approximate-svp-assumption}, we say that approximate CVP
assumption with the approximate factor $\gamma$ holds if
$\mathsf{Adv}^{\text{CVP}_{\gamma}}_{\mathcal{A},\mathsf{LGen}}(\mu)=\text{negl}(\mu)$
for all PTT adversaries $\mathcal{A}$.

\section{Learning With Errors (LWE)}\label{section:lwe}

\subsection{Key Generation}

Fix some modulo $q>1$ (further assume it is prime). Our primary goal is to build
public key cryptography schemes in the lattice-fashionable way. Turns out that
one of such approaches is to use \textbf{Learning With Errors} assumptions.

Fix secret key $\mathbf{s} \in \mathbb{Z}_q^n$ which we assume, for the moment,
is sampled uniformly randomly from the corresponding set (we will assume so
until \Cref{section:kyber}). The public key, corresponding to $\mathbf{s}$, will
be a vector of perturbed random inner products with $\mathbf{s}$. Specifically,
the public information consists of multiple pairs $(\mathbf{a}, b) \in
\mathbb{Z}_q^{n} \times \mathbb{Z}_q$ such that $b =
\langle\mathbf{a},\mathbf{s}\rangle + e$ for an \textit{error} $e \in
\mathbb{Z}_q$, typically taken from the distribution that produces very small
values of $e$: for instance, one might think of $e \in [-\eta,\eta]\pmod{q}$ for
$\eta \approx 3$. Denote such distribution by $\chi$.

This way, we can give an intermediate version of our key generation as follows:

\begin{center}
\begin{varwidth}{0.7\linewidth}
\begin{mdframed}[linecolor=black,
                 linewidth=1.0pt,
                 skipabove=10pt,
                 skipbelow=10pt,
                 leftmargin=0pt,
                 rightmargin=0pt]
\colorbox{purple!10}{\underline{\textbf{$\mathsf{KeyGen}(1^{\lambda}) \to (\mathsf{pk},\mathsf{sk})$:}}}

\begin{enumerate}
    \item $\mathbf{s} \leftarrowS \mathbb{Z}_q^{\lambda}$; \hspace{5px}{\textcolor{gray!70!black}{\footnotesize\texttt{// Generate secret}}}
    \item \underline{\textbf{for}} $i \in [\lambda]$:
    \item \hspace{10px} $\mathbf{a}_i \leftarrowS \mathbb{Z}_q^{\lambda}$, $e \leftarrowS \chi$;
    \item \hspace{10px} $b_i \gets \langle\mathbf{a}_i,\mathbf{s}\rangle + e$;\hspace{5px}{\footnotesize\texttt{\textcolor{gray!70!black}{// Generate random pertubation}}}
    \item \textbf{return} $(\mathsf{pk}=\{(\mathbf{a}_i, b_i)\}_{i \in [\lambda]}, \mathsf{sk}=\mathbf{s})$.
\end{enumerate}
\end{mdframed}
\end{varwidth}
\end{center}

It is then natural to batch all pertubations into a single matrix $A$ and vector
$\mathbf{b}$.\footnote{We also note that $A$ might not necessarily be (and is
typically not) the square matrix, but just for the sake of exposition we for now
assume $A \in \mathbb{Z}_q^{\lambda \times \lambda}$.} This way, our updated key
generation scheme is in \Cref{fig:keygen} (where by $\chi^{\lambda}$ we denote
the distribution over $\mathbb{Z}_q^{\lambda}$ where each component is taken
from $\chi$).

\begin{figure}[ht]
\centering
\begin{varwidth}{0.4\linewidth}
\begin{mdframed}[linecolor=black,
                 linewidth=1.0pt,
                 skipabove=10pt,
                 skipbelow=10pt,
                 leftmargin=0pt,
                 rightmargin=0pt]
\colorbox{purple!10}{\underline{$\mathsf{KeyGen}(1^{\lambda}) \to (\mathsf{pk},\mathsf{sk})$:}}
\begin{enumerate}
    \item $(\mathbf{s}, A) \leftarrowS \mathbb{Z}_q^{\lambda} \times \mathbb{Z}_q^{\lambda \times \lambda}$;
    \item $\mathbf{e} \leftarrowS \chi^{\lambda}$;
    \item $b \gets A\mathbf{s} + \mathbf{e}$;
    \item \textbf{return} $(\mathsf{pk}=(A,b), \mathsf{sk}=\mathbf{s})$.
\end{enumerate}
\end{mdframed}
\end{varwidth}
\caption{LWE-based Key Generation Scheme}
\label{fig:keygen}
\end{figure}

Note that we of course expect that solving $A\mathbf{s}+\mathbf{e}=\mathbf{b}$
is hard with respect to $\mathbf{s}$. It should also be clear why we introduce
the error term $\mathbf{e}$: assuming $\mathbf{e}=\mathbf{0}$, the equation
simplifies to $A\mathbf{s}=\mathbf{b}$. This is a simple system of linear
equations which can be further dealt with by means of, say, Gaussian
elimination, which has $O(\lambda^3)$ complexity. The assumption of hardness of
solving this equation for $\mathbf{e} \neq \mathbf{0}$ is what we introduce and
treat in the next sections.

\subsection{Search-LWE (SLWE) and Decisional-LWE (DLWE) Assumptions}

In this section we shall finally formalize main assumptions about learning with
errors problem. First, introduce the LWE instance generation algorithm as
specified in~\Cref{fig:lwe-instance} below.

\begin{figure}[ht]
\centering
\begin{varwidth}{0.45\linewidth}
\begin{mdframed}[linecolor=black,
                 linewidth=1.0pt,
                 skipabove=10pt,
                 skipbelow=10pt,
                 leftmargin=0pt,
                 rightmargin=0pt]
\colorbox{purple!10}{\underline{$\mathsf{LWEGen}_{n,m,q,\chi}(1^{\lambda}, \mathbf{s}) \to \mathbb{Z}_q^{m \times n} \times \mathbb{Z}_q^m$:}}
\begin{enumerate}
    \item $A \leftarrowS \mathbb{Z}_q^{m \times n}$;
    \item $\mathbf{e} \leftarrowS \chi$ (over $\mathbb{Z}_q^m$);
    \item $\mathbf{b} \gets A\mathbf{s} + \mathbf{e}$;
    \item \textbf{return} $(A, \mathbf{b})$.
\end{enumerate}
\end{mdframed}
\end{varwidth}
\caption{LWE Instance Generation}
\label{fig:lwe-instance}
\end{figure}

Now, informally, LWE assumption states that, given output of
$\mathsf{LWEGen}_{n,m,q,\chi}(1^{\lambda},\mathbf{s})$, it is impossible to
determine $\mathbf{s}$ (\textit{Search-LWE assumption}). Similarly to many
assumptions in other cryptographic systems (such as Diffie-Hellman assumption),
LWE also has the decisional analog: distinguish output of
$\mathsf{LWEGen}_{n,m,q,\chi}(1^{\lambda},\mathbf{s})$ from the uniform sample
over $\mathbb{Z}_q^{m \times n} \times \mathbb{Z}_q^m$ (\textit{Decisional-LWE
assumption}).  We state both assumptions more formally below.

\begin{definition}[Search-LWE Assumption]\label{def:search-lwe-assumption}
    For the given adversary $\mathcal{A}$ and LWE parameters $(n,m,q,\chi)$,
    define $\underline{\mathsf{Game} \;
    \mathsf{SLWE}}_{\mathsf{LWEGen}}^{\mathcal{A}}$ as specified
    in~\Cref{fig:search-lwe}. Define the advantage of $\mathcal{A}$ as
    \begin{equation*}
        \mathsf{Adv}^{\text{SLWE}}_{\mathcal{A},\mathsf{LWEGen}}(\lambda) = \text{Pr}\left[\underline{\mathsf{Game} \;
    \mathsf{SLWE}}_{\mathsf{LWEGen}}^{\mathcal{A}}(\lambda) = 1\right].
    \end{equation*}

    We say that \textbf{Search-LWE assumption} (SLWE) holds if
    $\mathsf{Adv}^{\text{SLWE}}_{\mathcal{A},\mathsf{LWEGen}}(\lambda)=\text{negl}(\lambda)$.
\end{definition}

\begin{definition}[Decisional-LWE Assumption]\label{def:decision-lwe-assumption}
    For the given adversary $\mathcal{A}$ and LWE parameters $(n,m,q,\chi)$,
    define $\underline{\mathsf{Game} \;
    \mathsf{DLWE}}_{\mathsf{LWEGen}}^{\mathcal{A}}$ as specified
    in~\Cref{fig:decision-lwe}. Define the advantage of $\mathcal{A}$ as
    \begin{equation*}
        \mathsf{Adv}^{\text{DLWE}}_{\mathcal{A},\mathsf{LWEGen}}(\lambda) = \left|\text{Pr}\left[\underline{\mathsf{Game} \;
    \mathsf{DLWE}}_{\mathsf{LWEGen}}^{\mathcal{A}}(\lambda) = 1\right] - \frac{1}{2}\right|.
    \end{equation*}

    We say that \textbf{Decisional-LWE assumption} (DLWE) holds if
    $\mathsf{Adv}^{\text{DLWE}}_{\mathcal{A},\mathsf{LWEGen}}(\lambda)=\text{negl}(\lambda)$
    for all PPT adversaries $\mathcal{A}$.
\end{definition}

\begin{figure}[ht]
    \centering
    \begin{minipage}{0.45\linewidth}
        \begin{mdframed}[linecolor=black, linewidth=1pt,
                         skipabove=10pt, skipbelow=10pt]
        \colorbox{green!10}{\underline{$\mathsf{Game} \; \mathsf{SLWE}_{\mathsf{LWEGen}}^{\mathcal{A}}$:}}
        \begin{enumerate}
            \item $\mathbf{s} \leftarrowS \mathbb{Z}_q^{n}$;
            \item $(A,\mathbf{b}) \leftarrowS \mathsf{LWEGen}_{n,m,q,\chi}(\mathbf{s})$;
            \item $\widehat{\mathbf{s}} \gets \mathcal{A}(A,\mathbf{b})$;
            \item \textbf{return} $\mathbf{s} =_? \widehat{\mathbf{s}}$.
        \end{enumerate}
        \end{mdframed}
        \caption{Search-LWE Game} 
        \label{fig:search-lwe}
    \end{minipage}
    \hfill
    \begin{minipage}{0.45\linewidth}
        \begin{mdframed}[linecolor=black, linewidth=1pt,
                         skipabove=10pt, skipbelow=10pt]
        \colorbox{green!10}{\underline{$\mathsf{Game} \; \mathsf{DLWE}_{\mathsf{LWEGen}}^{\mathcal{A}}$:}}
        \begin{enumerate}
            \item $\mathbf{s} \leftarrowS \mathbb{Z}_q^{n}$;
            \item $(A_0,\mathbf{b}_0) \leftarrowS \mathsf{LWEGen}_{n,m,q,\chi}(\mathbf{s})$;
            \item $(A_1,\mathbf{b}_1) \leftarrowS \mathbb{Z}_q^{m \times n} \times \mathbb{Z}_q^m$;
            \item $c \leftarrowS \{0,1\}$;
            \item $\widehat{c} \gets \mathcal{A}(A_c,\mathbf{b}_c)$;
            \item \textbf{return} $c =_? \widehat{c}$.
        \end{enumerate}
        \end{mdframed}
        \caption{Decision-LWE Game}
        \label{fig:decision-lwe}
    \end{minipage}
\end{figure}

The natural question to ask is whether SLWE and DLWE assumptions are equivalent
or not. In fact, \cite{regev} has proven that these notions are in fact
equivalent in case $q=\mathsf{poly}(n)$. We show this in the following
proposition.

\begin{proposition}[Following~\cite{regev}]
    If $q=\mathsf{poly}(n)$, SLWE and DLWE assumptions are polynomially
    equivalent. To put it the other way, if the SLWE assumption holds for LWE
    instance with parameters $(n,m,q=\mathsf{poly}(n),\chi)$, then the DLWE
    assumption also holds.
\end{proposition}

\textbf{Proof.} $\text{SLWE} \Rightarrow \text{DLWE}$ reduction is obvious. We
therefore show $\text{DLWE} \Rightarrow \text{SLWE}$. Suppose exists adversary
$\mathcal{A}$ that wins $\underline{\mathsf{Game} \;
\mathsf{DLWE}}_{\mathsf{LWEGen}}^{\mathcal{A}}$ in a non-negligible probability.
We construct $\mathcal{B}^{\mathcal{A}}$ that wins $\underline{\mathsf{Game} \;
\mathsf{SLWE}}_{\mathsf{LWEGen}}^{\mathcal{B}}$ is a non-negligible probability
as well.

The informal idea is as follows: suppose $\mathcal{B}$ wants to solve
$A\mathbf{s}+\mathbf{e}=\mathbf{b}$. Turns out that $\mathcal{B}$ can guess the
coordinates of $\mathbf{s}$ one at a time, starting from $s_1$. Suppose
$\mathcal{B}$ picks some $\widetilde{s} \in \mathbb{Z}_q$ and wants to check
whether $\widetilde{s} = s_1$. Select some $\boldsymbol{\Delta} \leftarrowS
\mathbb{Z}_q^m$ and let
$\widetilde{A}=(\widetilde{\mathbf{a}}_1,\mathbf{a}_2,\dots,\mathbf{a}_n)$ where
$\widetilde{\mathbf{a}}_1=\mathbf{a}_1+\boldsymbol{\Delta}$ (and where the
initial matrix is given by $A=(\mathbf{a}_1,\dots,\mathbf{a}_n)$) and set
$\widetilde{\mathbf{b}} := \mathbf{b} + \widetilde{s}\boldsymbol{\Delta}$. Now,
if $\widetilde{s}=s_1$, then $\widetilde{\mathbf{b}} = \widetilde{A}\mathbf{s} +
\mathbf{e}$, otherwise $\widetilde{\mathbf{b}} = \widetilde{A}\mathbf{s} +
\mathbf{e}+(\widetilde{s}-s_1)\boldsymbol{\Delta}$. Since $\boldsymbol{\Delta}$
is uniform, it follows that in the latter case $\widetilde{\mathbf{b}}$ is
uniformly random and independent of $\widetilde{A}$. Thus, using $\mathcal{A}$,
adversary $\mathcal{B}$ has an effective test for chosen $\widetilde{s}$, and
thus $\mathcal{B}$ can range over all possible $q$ values, which is polynomial
in $n$.

This way, by using $\mathcal{B}^{\mathcal{A}}_{\text{SLWE}}$ defined in
\Cref{fig:dlwe-to-slwe}, we obtain:
\begin{equation*}
    \mathsf{Adv}^{\text{DLWE}}_{\mathcal{A},\mathsf{LWEGen}}(\lambda) \leq qn \cdot \mathsf{Adv}^{\text{SLWE}}_{\mathcal{B}^{\mathcal{A}},\mathsf{LWEGen}}(\lambda). \quad \qedwhite
\end{equation*}

\begin{remark}
    Note that it might seem like the factor of $qn$ is too large to state the
    equivalence of DLWE and SLWE (indeed, in the elliptic curve cryptography,
    one typically uses $q$ exponential in the security parameter). However, in practical
    applications, $q$ is typically not considerably large: for instance, in
    Kyber~\cite{kyber}, $q=7681$ and $n=256$!

    Also note that the brute-force complexity without the DLWE oracle is $q^n$,
    which is of course overwhelmingly large! (compared to $qn$).
\end{remark}

\begin{figure}[ht]
\centering
\begin{varwidth}{0.7\linewidth}
\begin{mdframed}[linecolor=black,
                 linewidth=1.0pt,
                 skipabove=10pt,
                 skipbelow=10pt,
                 leftmargin=0pt,
                 rightmargin=0pt]
\colorbox{purple!10}{\underline{$\mathcal{B}^{\mathcal{A}}_{\text{SLWE}}(1^{\lambda}, A=(\mathbf{a}_1,\dots,\mathbf{a}_n), \mathbf{b}) \to \mathbb{Z}_q^{m}$:}}
\begin{enumerate}
    \item $\widehat{\mathbf{s}} \gets \mathbf{0} \in \mathbb{Z}_q^m$;\hspace{5px}{\textcolor{gray!70!black}{\footnotesize\texttt{// Prepare the output}}}
    \item \textbf{for} $i \in [n]$:
    \item \hspace{10px} \textbf{for} $\widetilde{s} \in \mathbb{Z}_q$:
    \item \hspace{20px} $\boldsymbol{\Delta} \leftarrowS \mathbb{Z}_q^m$;\hspace{5px}{\textcolor{gray!70!black}{\footnotesize\texttt{// Generate random offset}}}
    \item \hspace{20px} $\widetilde{A} \gets (\mathbf{a}_1,\dots,\mathbf{a}_{i-1},\widetilde{\mathbf{a}}_i, \mathbf{a}_{i+1},\dots,\mathbf{a}_n)$, $\widetilde{\mathbf{a}}_i := \mathbf{a}_i+\boldsymbol{\Delta}$;
    \item \hspace{20px} $\widetilde{\mathbf{b}} \gets \mathbf{b} + \widetilde{s}\boldsymbol{\Delta}$;\hspace{5px}{\textcolor{gray!70!black}{\footnotesize\texttt{// Set offset vector $\widetilde{\mathbf{b}}$ and matrix $\widetilde{A}$}}}
    \item \hspace{20px} $c \gets \mathcal{A}(\widetilde{A}, \widetilde{\mathbf{b}})$;\hspace{5px}{\textcolor{gray!70!black}{\footnotesize\texttt{// Run DLWE oracle}}}
    \item \hspace{20px} \textbf{if} $c=0$, set $\widehat{s}_i \gets \widetilde{s}$ and \textbf{break};
    \item \textbf{return} $\widehat{\mathbf{s}}$.
\end{enumerate}
\end{mdframed}
\end{varwidth}
\caption{Algorithm for finding the solution to LWE instance given oracle to the DLWE.}
\label{fig:dlwe-to-slwe}
\end{figure}

\subsection{LWE-based Public-Key Encryption}

Finally, let us define the public-key encryption (PKE) scheme based on SLWE and
DLWE assumptions! We remind the reader of the definition of a PKE scheme.
\begin{definition}[Public-Key Encryption Scheme]
    A \textbf{public-key encryption scheme} $\Pi_{\text{PKE}}$ is a triplet 
    of the following PPT algorithms:
    \begin{itemize}
        \item $\mathsf{KeyGen}(1^{\lambda}) \to (\mathsf{pk},\mathsf{sk})$
        returns the pair of public key and secret key, respectively.
        \item $\mathsf{Enc}(\mathsf{pk}, m) \to c$ takes the public key and message
        and produces the ciphertext $c$.
        \item $\mathsf{Dec}(\mathsf{sk}, c) \to m$ takes the secret key and ciphertext and restores the message $m$.
    \end{itemize}
    We require \textbf{correctness}:
    \begin{equation*}
        \forall m: \; \text{Pr}\left[ \widehat{m} = m \;\middle|\; \begin{aligned}
            &(\mathsf{pk}, \mathsf{sk}) \gets \mathsf{KeyGen}(1^{\lambda}) \\
            &\widehat{m} \gets \mathsf{Dec}(\mathsf{sk}, \mathsf{Enc}(\mathsf{pk}, m))
        \end{aligned} \right] \geq 1 - \text{negl}(\lambda),
    \end{equation*}
    and \textbf{IND-CCA} and \textbf{IND-CPA} security. To not overwhelm the reader 
    with the details, we will only focus on the IND-CPA security with the corresponding $\underline{\mathsf{Game} \; \mathsf{INDCPA}}^{\mathcal{A}}_{\Pi}$ defined in \Cref{fig:ind-cpa-game}. We 
    naturally define advantage $\mathsf{Adv}_{\mathcal{A},\Pi_{\text{PKE}}}^{\text{IND-CPA}}(\lambda) := 2\left|\text{Pr}[\underline{\mathsf{Game} \; \mathsf{INDCPA}}^{\mathcal{A}}_{\Pi_{\text{PKE}}} = 1] - \frac{1}{2}\right|$ and say that $\Pi_{\text{PKE}}$ is \textbf{IND-CPA secure} if $\mathsf{Adv}_{\mathcal{A},\Pi_{\text{PKE}}}^{\text{IND-CPA}}(\lambda) = \text{negl}(\lambda)$ for all PPT $\mathcal{A}$.
\end{definition}

\begin{figure}[ht]
\centering
\begin{varwidth}{0.75\linewidth}
\begin{mdframed}[linecolor=black, linewidth=1pt,
                    skipabove=10pt, skipbelow=10pt]
\colorbox{green!10}{\underline{$\mathsf{Game} \; \mathsf{INDCPA}_{\Pi}^{\mathcal{A}}(1^{\lambda})$:}}
    \begin{enumerate}
        \item $(\mathsf{pk},\mathsf{sk}) \gets \mathsf{KeyGen}(1^{\lambda})$;
        \item $(m_0,m_1) \gets \mathcal{A}(\mathsf{pk})$;\hspace{5px}{\textcolor{gray!70!black}{\footnotesize\texttt{// Typically one requires $|m_0|=|m_1|$}}}
        \item $b \leftarrowS \{0,1\}$, $c_b \gets \mathsf{Enc}(\mathsf{pk}, m_b)$;
        \item $\widehat{b} \gets \mathcal{A}(\mathsf{pk}, m_0, m_1, c_b)$;
        \item \textbf{return} $\widehat{b} =_? b$.
    \end{enumerate}
\end{mdframed}
\end{varwidth}
\caption{IND-CPA Security Game. $\mathcal{A}$ might be
allowed to perform steps 2--3 $\text{poly}(\lambda)$ times.} 
\label{fig:ind-cpa-game}
\end{figure}

Now, we define the following PKE scheme based on $(n,m,q,\chi)$ assumption where
the message space is $\{0,1\}$. The main idea is that encryption scheme will
generate the new LWE instance $(A',\mathbf{b}')$ by taking a ``small'' linear
combination of publicly known $(A,\mathbf{b})$. By additionally adding
$m\lceil q/2 \rfloor$ to $\mathbf{b}'$, we obtain the encryption (further denote
by $\lceil \cdot \rfloor: \mathbb{Q} \to \mathbb{Z}$ the function that rounds a
given rational number to the nearest integer). Upon decryption, the holder of 
the secret key $\mathbf{s}$ will be able to check whether $A'\mathbf{s}-\mathbf{b}'$
produces the small error. If the term $m\lceil q/2 \rfloor$ was added to form 
$\mathbf{b}'$, the resultant difference would be large, suggesting that the original 
bit was $1$. Otherwise, decryption would output $0$. Let us put this formally in 
the definition below.

\begin{definition}[LWE-based PKE]\label{def:lwe-pke}
    Based on LWE instance with parameters $(n,m,q,\chi)$, define the PKE scheme
    $\Pi_{\text{PKE}}^{\text{LWE}}$ over message space $\{0,1\}$ as follows: $\mathsf{KeyGen}$ is specified in
    \Cref{fig:keygen}, $\mathsf{Enc}$ in \Cref{fig:lwe-enc}, and $\mathsf{Dec}$
    in \Cref{fig:lwe-dec}.
\end{definition}

Since we have to work with ``small'' and ``large'' elements of $\mathbb{Z}_q$,
by $\|x\|_{\infty}$ for $x \in \mathbb{Z}_q$ we denote $x$ if $x \leq \frac{q}{2}$, and
$q-x$ otherwise (for example, naturally, $\|q-3\|_{\infty} = 3$).

\begin{figure}[ht]
    \centering
    \begin{minipage}{0.45\linewidth}
        \begin{mdframed}[linecolor=black, linewidth=1pt,
                         skipabove=10pt, skipbelow=10pt]
        \colorbox{green!10}{\underline{$\mathsf{Enc}(\mathsf{pk}, m \in \{0,1\}) \to (c \in \mathbb{Z}_q^n \times \mathbb{Z}_q)$:}}
        \begin{enumerate}
            \item $\Delta\mathbf{r}, \Delta\mathbf{z} \leftarrowS \chi^m$, $\Delta z' \leftarrowS \chi$;
            \item $\widetilde{\mathbf{a}} \gets A^{\top}\Delta\mathbf{r} + \Delta\mathbf{z}$;
            \item $\widetilde{b} \gets \mathbf{b}^{\top}\Delta\mathbf{r}  + \Delta z' + m\lceil q/2 \rfloor$;
            \item \textbf{return} $c=(\widetilde{\mathbf{a}}, \widetilde{b})$.
        \end{enumerate}
        \end{mdframed}
        \caption{Encryption in $\Pi_{\text{PKE}}^{\text{LWE}}$} 
        \label{fig:lwe-enc}
    \end{minipage}
    \hfill
    \begin{minipage}{0.45\linewidth}
        \begin{mdframed}[linecolor=black, linewidth=1pt,
                         skipabove=10pt, skipbelow=10pt]
        \colorbox{green!10}{\underline{$\mathsf{Dec}(\mathsf{sk}, c) \to m$:}}
        \begin{enumerate}
            \item Parse $(\widetilde{\mathbf{a}}, \widetilde{b}) \gets c$;
            \item $\delta \gets \widetilde{b} - \mathbf{s}^{\top}\widetilde{\mathbf{a}}$;
            \item \textbf{return} $1$ iff $\|\delta\|_{\infty} > q/4$.
        \end{enumerate}
        \end{mdframed}
        \caption{Decryption in $\Pi_{\text{PKE}}^{\text{LWE}}$}
        \label{fig:lwe-dec}
    \end{minipage}
\end{figure}

Now we show that $\Pi_{\text{PKE}}^{\text{LWE}}$ is a correct and secure PKE scheme.

\begin{proposition}\label{prop:ind-cpa-security}
    $\Pi_{\text{PKE}}^{\text{LWE}}$ is correct and IND-CPA secure under DLWE
    assumption.
\end{proposition}

\textbf{Proof.} (\textit{Correctness}) Let us first show correctness. Assume $\mathsf{sk}=\mathbf{s}$
and $\mathsf{pk}=(A,\mathbf{b})$ is the valid key-pair (that is, $A\mathbf{s}+\mathbf{e}=\mathbf{b}$ for ``small'' $\mathbf{e} \leftarrowS \chi^m$) and we formed
$(\widetilde{\mathbf{a}}, \widetilde{b})$ according to \Cref{fig:lwe-enc}. Let 
us try to decode this cipher by means of \Cref{fig:lwe-dec}:
\begin{align*}
    \delta &= \mathbf{b}^{\top}\Delta\mathbf{r} + \Delta z' + m\lceil q/2 \rfloor - 
    \mathbf{s}^{\top}A^{\top}\Delta\mathbf{r} - \mathbf{s}^{\top}\Delta\mathbf{z} \\
    &= (\mathbf{b} - A\mathbf{s})^{\top}\Delta \mathbf{r} + (\Delta z' - \mathbf{s}^{\top}\Delta\mathbf{z}) + m\lceil q/2 \rfloor \\
    &= \textcolor{blue!70!black}{(\mathbf{e}^{\top}\Delta \mathbf{r} - \mathbf{s}^{\top}\Delta\mathbf{z} + \Delta z')} + m\lceil q/2 \rfloor
\end{align*}

Now note that the quantity marked in \textcolor{blue!70!black}{blue} is small.
Indeed, let us show why. Assume that $\chi$ produces elements of $\mathbb{Z}_q$
with the modulus not exceeding $\eta$. Then\footnote{In the sequel, we will show
that $\mathbf{s}$ can be taken from the same distribution as the error term
$\mathbf{e}$ (that is, from $\chi$), making upper bound even smaller.},
\begin{equation*}
    \|\textcolor{blue!70!black}{\mathbf{e}^{\top}\Delta \mathbf{r} - \mathbf{s}^{\top}\Delta\mathbf{z} + \Delta z'}\|_{\infty} \leq \|\mathbf{e}^{\top}\Delta \mathbf{r}\|_{\infty} + \|\mathbf{s}^{\top}\Delta\mathbf{z}\|_{\infty} + \|\Delta z'\|_{\infty} \leq n\eta^2 + n\eta + \eta \leq 2n\eta.
\end{equation*}

Clearly, $2n\eta$ is a small quantity (assuming $\eta \ll q/8n$). Now, if $m=0$, we
get that $\|\delta\|_{\infty} = \|(\textcolor{blue!70!black}{\text{blue part}})\|_{\infty} < q/4$, so
decryption works. If $m=1$, then $\delta = \lceil q/2\rfloor +
(\textcolor{blue!70!black}{\text{blue part}})$, which is a large value, 
so the decryption outputs $1$, as expected.

(\textit{IND-CPA Security}). Suppose $\mathcal{A}$ solves the
$\underline{\mathsf{Game} \; \mathsf{INDCPA}}_{\Pi}^{\mathcal{A}}$ with a
non-negligible probability. We construct the $\mathcal{B}^{\mathcal{A}}$ that
breaks the $\underline{\mathsf{Game} \;
\mathsf{DLWE}}_{\mathsf{LWEGen}}^{\mathcal{B}}$ using $\mathcal{A}$ as an
oracle. The high-level idea of the proof is the following: suppose $\mathcal{B}$
is given a tuple $(A,\mathbf{b})$. If $(A,\mathbf{b})$ is a valid LWE instance,
then it is a valid public key in $\Pi_{\text{PKE}}^{\text{LWE}}$, so
$\mathcal{B}$ might run $\mathcal{A}$ on $\mathsf{pk}^* \gets (A,\mathbf{b})$.
Then, if $\mathcal{A}$ correctly guesses the bit, then likely $\mathsf{pk}^*$
was a valid key. In turn, if $\mathsf{pk}^*$ is a ``garbage'', with probability
approximately $1/2$, the $\mathcal{A}$ would lose. Thus, $\mathcal{B}$ decision
will be based on whether $\mathcal{A}$ guesses the correct bit or not. We
encapsulate this in \Cref{fig:ind-cpa-to-dlwe}. One can then show:
\begin{equation*}
    \mathsf{Adv}^{\text{IND-CPA}}_{\mathcal{A},\Pi}(\lambda) \leq 2\mathsf{Adv}^{\text{DLWE}}_{\mathcal{B}^{\mathcal{A}},\mathsf{LWEGen}}(\lambda). 
    \quad \qedwhite
\end{equation*}

\begin{figure}[ht]
\centering
\begin{varwidth}{0.7\linewidth}
\begin{mdframed}[linecolor=black,
                 linewidth=1.0pt,
                 skipabove=10pt,
                 skipbelow=10pt,
                 leftmargin=0pt,
                 rightmargin=0pt]
\colorbox{purple!10}{\underline{$\mathcal{B}^{\mathcal{A}}_{\text{DLWE}}(1^{\lambda}, A, \mathbf{b}) \to \{0,1\}$:}}
\begin{enumerate}
    \item pass $\mathsf{pk}^* \gets (A,\mathbf{b})$ to $\mathcal{A}$;
    \item get $(m_0,m_1)$ from $\mathcal{A}$; 
    \item choose $d \leftarrowS \{0,1\}$;
    \item send $c_d \gets \mathsf{Enc}(\mathsf{pk}^*, m_d)$ to $\mathcal{A}$;
    \item $\mathcal{A}$ responds with $\widehat{d}$;
    \item \textbf{return} $d =_? \widehat{d}$. 
\end{enumerate}
\end{mdframed}
\end{varwidth}
\caption{Algorithm for finding the solution to DLWE instance given oracle to the IND-CPA.}
\label{fig:ind-cpa-to-dlwe}
\end{figure}

\subsection{Relation to Lattices}

The reasonable question that the reader might rightfully have at this point:
where are the lattices? So far, we have spent the whole \Cref{section:lattices}
describing lattices, while the current section does not explicitly use them in
any way.

In fact, turns out that the LWE problem can be viewed as the lattice problem!

\begin{definition}
    For the LWE instance $(A,\mathbf{b})$ over parameters $(n,m,q,\chi)$, we
    define the associated lattice $\Lambda_q(A)$ as follows:
    \begin{equation*}
        \Lambda_q(A) \triangleq \{\mathbf{y} \in \mathbb{Z}^m: \mathbf{y} = A\mathbf{x} \; (\text{mod}\; q) \; \text{for some} \; \mathbf{x} \in \mathbb{Z}^n\} = A\mathbb{Z}^n + q\mathbb{Z}^m.
    \end{equation*}
\end{definition}

Now note that the vector $\mathbf{b}$ does not lie on the lattice $\Lambda_q(A)$
(unless $\mathbf{e}=\mathbf{0}$ which makes LWE instance unsecure). But note
that due to the modular equation, $\mathbf{b} = A\mathbf{s} + q\mathbf{z} +
\mathbf{e}$ for some $\mathbf{z} \in \mathbb{Z}^m$ and since $A\mathbf{s} +
q\mathbf{z} \in \Lambda_q(A)$, the distance from $\mathbf{b}$ to $\Lambda_q(A)$
is bounded by $\|\mathbf{e}\|$. That said, although $\mathbf{b}$ is not in the
lattice, it is very close to it! Therefore, almost surely, the closest vector to
$\mathbf{b}$ is $A\mathbf{s} + q\mathbf{z}$. Now imagine that we know how to
solve CVP defined in \Cref{section:approximate-cvp}. Then, based on
$\mathbf{b}$, we can derive $A\mathbf{s}+q\mathbf{z}$, and then easily obtain
$\mathbf{s}$ by solving a system of linear equations. Note that the
approximation factor depends on the underlying distribution $\chi^m$, from which
$\mathbf{e}$ is sampled. 

This way, lattices are crucial for cryptanalysis of LWE-based schemes.

\subsection{Variations of LWE}

Let us analyze the efficiency of the PKE scheme provided in \Cref{def:lwe-pke}.
\begin{proposition}
    In the PKE scheme $\Pi_{\text{PKE}}^{\text{LWE}}$ over LWE instance
    $(n,m,q,\chi)$, the public key size is $O(m\log q)$, secret key size is
    $O(n\log q)$, and the codeword size is $O(m\log q)$: therefore, sending a
    $B$-bit message takes $O(B m \log q)$ space.
\end{proposition}

\begin{remark}\label{remark:seeded-a}
    One might ask why the public key size is not $O(nm\log q)$, given that we need
to store the matrix $A$. Fortunately, $A$ is sampled randomly, so instead of
storing the whole matrix, one might store the seed from which one can derive
$A$ instead.
\end{remark}

Unfortunately, such key and ciphertext sizes are significantly larger than those
of elliptic curve-based schemes (in fact, if $\lambda$ is taken as a security
parameter, then ciphertext size scales as $O(\lambda^2)$). This led to the
natural question: can we reduce these sizes?

\subsubsection{Ring-LWE}

Turns out that one possible answer to this question is to use objects other than
$\mathbb{Z}_q^n$ as elements that possess certain useful structure: for
instance, the popular choice is to use $\mathcal{R}_q :=
\mathcal{R}/q\mathcal{R}$ for $\mathcal{R} = \mathbb{Z}[T]/(T^n+1)$ and $n$
being a power-of-two. In such case, one might think of $\mathcal{R}$ either by: 
\begin{enumerate}
    \item[(a)] $(n-1)$-degree polynomials with operations modulo $T^n+1$;
    \item[(b)] the ring of integers $\mathbb{Z}(\zeta_{2n})$ since $\mathcal{R}
\cong \mathbb{Z}[T]/(\Phi_{2n}(T)) \cong \mathbb{Z}(\zeta_{2n})$ where $\zeta_N
:= e^{2\pi i/N}$.
\end{enumerate}

Now, instead of having $(A,\mathbf{b}) \in \mathbb{Z}_q^{m \times n} \times
\mathbb{Z}_q^m$ as a public key, we will store the public key as a pair $(a,b)
\in \mathcal{R}_q^2$ where $b=as+e$, but over the polynomial space. However, we
should clarify what ``small'' means in the context of $\mathcal{R}_q$ (since $e$
here comes from some small distribution). For this reason, given $f(T)=\sum_{i
\in (n)}f_iT^i \in \mathcal{R}_q$, define:
\begin{equation*}
    \|f\|_{\infty} \triangleq \max_{i \in (n)}\|f_i\|_{\infty}.
\end{equation*}
So in other words, polynomial $f(T) \in \mathcal{R}_q$ is large as long as its
largest coefficient is large. 

Further, we identify Ring-LWE instance by a tuple of $(\mathcal{R},n,q,\chi)$,
where (\textit{i}) additive group of $\mathcal{R}$ is isomorphic to
$\mathbb{Z}^n$, (\textit{ii}) instead of $\mathbb{Z}_q^n$ in ``plain'' LWE, we
use field $\mathcal{R}_q := \mathcal{R}/q\mathcal{R}$, and (\textit{iii}) $\chi$
is a (small) random distribution over $\mathcal{R}_q$.

\begin{proposition}
    Using Ring-LWE structure, we can reduce ciphertext size for a $n$-bit message
    of PKE Scheme from \Cref{def:lwe-pke} from $O(nm\log q)$ to $O(n\log q)$.
    Moreover, computing expression $as+e$ now takes only $O(n\log n)$ complexity
    instead of $O(nm)$.
\end{proposition}

The reason for $O(n\log n)$ complexity for computing $as+e$ comes from the fact
that we can utilize Number Theoretic Transform (NTT/FFT) for fast polynomial
multiplication. The reason for reduced ciphertext size is less evident, but the
idea is the following: instead of adding a $m\lceil q/2 \rfloor$ term to the
second ciphertext piece in \Cref{fig:lwe-enc}, we are going to represent $n$
bits of message $(m_i)_{i \in (n)}$ as a polynomial $m(T) := \sum_{i \in
(n)}m_iT^i \in \mathcal{R}_q$. We specify details in the definition below.

\begin{definition}\label{def:rlwe-pke}
    Based on Ring-LWE instance with parameters $(\mathcal{R}, n, q,\chi)$, define the PKE scheme
    $\Pi_{\text{PKE}}^{\text{RLWE}}$ over message space $\{0,1\}^n$ as follows: $\mathsf{KeyGen}$ is specified in
    \Cref{fig:ring-keygen}, $\mathsf{Enc}$ in \Cref{fig:ring-lwe-enc}, and $\mathsf{Dec}$
    in \Cref{fig:ring-lwe-dec}.
\end{definition}

\begin{figure}[ht]
\centering
\footnotesize 

\begin{minipage}[t]{0.32\linewidth}
    \begin{mdframed}[linecolor=black, linewidth=1.0pt,
                     skipabove=5pt, skipbelow=5pt,
                     innerleftmargin=5pt, innerrightmargin=5pt]
        \centering
        \colorbox{purple!10}{\underline{$\mathsf{KeyGen}(1^{\lambda}) \to (\mathsf{pk},\mathsf{sk})$:}}
        \begin{enumerate}
            \item $(s,a) \leftarrowS \mathcal{R}_q^2$;
            \item $e \leftarrowS \chi$ (over $\mathcal{R}_q$);
            \item $b \gets as + e \in \mathcal{R}_q$;
            \item \textbf{return} $((a,b),s)$.
        \end{enumerate}
    \end{mdframed}
    \caption{Ring-LWE-based Key Generation}
    \label{fig:ring-keygen}
\end{minipage}
\hfill
\begin{minipage}[t]{0.32\linewidth}
    \begin{mdframed}[linecolor=black, linewidth=1.0pt,
                     skipabove=5pt, skipbelow=5pt,
                     innerleftmargin=5pt, innerrightmargin=5pt]
        \centering
        \colorbox{green!10}{\underline{$\mathsf{Enc}(\mathsf{pk}, (m_i)_{i \in (n)} \in \{0,1\}^n) \to c$:}}
        \begin{enumerate}
            \item $m \gets \sum_{i \in (n)}m_iT^i \in \mathcal{R}_q$;
            \item $\Delta r, \Delta z, \Delta z' \leftarrowS \chi$;
            \item $\widetilde{a} \gets a\Delta r + \Delta z$;
            \item $\widetilde{b} \gets b\Delta r  + \Delta z' + \lceil \frac{q}{2} \rfloor m$;
            \item \textbf{return} $c=(\widetilde{a}, \widetilde{b})$.
        \end{enumerate}
    \end{mdframed}
    \caption{Ring-LWE-based Encryption} 
    \label{fig:ring-lwe-enc}
\end{minipage}
\hfill
\begin{minipage}[t]{0.32\linewidth}
    \begin{mdframed}[linecolor=black, linewidth=1.0pt,
                     skipabove=5pt, skipbelow=5pt,
                     innerleftmargin=5pt, innerrightmargin=5pt]
        \centering
        \colorbox{green!10}{\underline{$\mathsf{Dec}(\mathsf{sk}, c) \to m$:}}
        \begin{enumerate}
            \item Parse $(\widetilde{a}, \widetilde{b}) \gets c$;
            \item $\delta \gets \widetilde{b} - s\widetilde{a}$;
            \item \textbf{return} $1$ iff $\|\delta\|_{\infty} > \frac{q}{4}$.
        \end{enumerate}
    \end{mdframed}
    \caption{Ring-LWE-based Decryption}
    \label{fig:ring-lwe-dec}
\end{minipage}
\end{figure}

Another valid question is whether introducing such structure should worsen the
security: so in other words, given appropriate Ring-LWE instance
$(\mathcal{R},n,q,\chi)$, is $\Pi_{\text{PKE}}^{\text{RLWE}}$ secure? Turns out
that there are no algorithms for breaking analogous Ring-LWE-based assumption
significantly faster than ``plain'' version of the LWE. To put this assumption
formally, we slighly modify \Cref{def:search-lwe-assumption}.

\begin{definition}
    For the given adversary $\mathcal{A}$ and Ring-LWE parameters
    $(\mathcal{R},n,q,\chi)$, define $\underline{\mathsf{Game} \;
    \mathsf{SRLWE}}^{\mathcal{A}}$ as specified in~\Cref{fig:search-rlwe}. Define
    the advantage of $\mathcal{A}$ as
    \begin{equation*}
        \mathsf{Adv}^{\text{SRLWE}}_{\mathcal{A}}(\lambda) = \text{Pr}\left[\underline{\mathsf{Game} \;
    \mathsf{SRLWE}}^{\mathcal{A}}(\lambda) = 1\right].
    \end{equation*}

    We say that \textbf{Search-Ring-LWE assumption} (SRLWE) holds if
    $\mathsf{Adv}^{\text{SRLWE}}_{\mathcal{A}}(\lambda)=\text{negl}(\lambda)$
    for all PPT adversaries $\mathcal{A}$. Decisional version of Ring-LWE is
    defined similarly to \Cref{def:decision-lwe-assumption}.
\end{definition}

\begin{figure}[ht]
\centering
\begin{varwidth}{0.4\linewidth}
\begin{mdframed}[linecolor=black, linewidth=1pt,
                         skipabove=10pt, skipbelow=10pt]
\colorbox{green!10}{\underline{$\mathsf{Game} \; \mathsf{SRLWE}^{\mathcal{A}}$:}}
\begin{enumerate}
    \item $(s,a) \leftarrowS \mathcal{R}_q^2$;
    \item $e \leftarrowS \chi$;
    \item $b \gets as + e$;
    \item $\widehat{s} \gets \mathcal{A}(a,b)$;
    \item \textbf{return} $s =_? \widehat{s}$.
\end{enumerate}
\end{mdframed}
\end{varwidth}
\caption{Search-Ring-LWE Game} 
\label{fig:search-rlwe}
\end{figure}

We leave the following proposition without a proof as it closely resembles one
specified in Proposition~\ref{prop:ind-cpa-security}. For concrete details and
survey on possible attacks, see for example \cite{lac}.
\begin{proposition}
    $\Pi_{\text{PKE}}^{\text{RLWE}}$ is correct and IND-CPA secure under Decisional Ring-LWE
    assumption.
\end{proposition}

\subsubsection{Module-LWE}

The idea of Module-LWE is to in a way ``blend'' together Ring-LWE and plain-LWE.
Namely, instead of using $\mathcal{R}_q$ in Ring-LWE, we will use vectors from
$\mathcal{R}_q^{\ell}$. Note that equipping $\mathcal{R}_q$ with a vector
structure naturally induces metric notions we have used before: e.g., for
$\bm{f}=(f_1,\dots,f_{\ell}) \in \mathcal{R}_q^{\ell}$, we have
$\|\bm{f}\|_{\infty} := \max_{i \in [\ell]}\|f_i\|_{\infty}$ and $\langle
\bm{f}, \bm{g} \rangle = \sum_{i \in [\ell]}f_ig_i$ is a valid inner product
that produces output in $\mathcal{R}_q$. 

The analogy to plain-LWE equation $\mathbf{b} = A\mathbf{s}+\mathbf{e}$ for $A
\in \mathbb{Z}_q^{n \times m}$, $\mathbf{s} \in \mathbb{Z}_q^m$, and
$\mathbf{b},\mathbf{e} \in \mathbb{Z}_q^n$, is now $\mathbf{b}_{\mathcal{R}} =
A_{\mathcal{R}}\mathbf{s}_{\mathcal{R}} + \mathbf{e}_{\mathcal{R}}$ for
$A_{\mathcal{R}} \in \mathcal{R}_q^{k \times \ell}$, $\mathbf{s} \in
\mathcal{R}_q^{\ell}$, and $\mathbf{s}_{\mathcal{R}},\mathbf{e}_{\mathcal{R}}
\in \mathcal{R}_q^{k}$, where operations are done inside ring $\mathcal{R}_q$.
We call $(A,\mathbf{b})$ a Module-LWE instance with parameters
$(\mathcal{R},k,\ell,q,\chi)$ if it is an output of $\mathsf{MLWEGen}_{\mathcal{R},k,\ell,q,\chi}$
defined in \Cref{fig:mlwe-gen}. Now we define Module-LWE assumption in the 
definition below.

\begin{definition}
    For the given adversary $\mathcal{A}$ and Module-LWE parameters
    $(\mathcal{R},k,\ell,q,\chi)$, define $\underline{\mathsf{Game} \;
    \mathsf{SMLWE}}^{\mathcal{A}}$ as specified in~\Cref{fig:mlwe-game}. Define
    the advantage of $\mathcal{A}$ as
    \begin{equation*}
        \mathsf{Adv}^{\text{SMLWE}}_{\mathcal{A}}(\lambda) = \text{Pr}\left[\underline{\mathsf{Game} \;
    \mathsf{SMLWE}}^{\mathcal{A}}(\lambda) = 1\right].
    \end{equation*}

    We say that \textbf{Search-Module-LWE assumption} (SMLWE) holds if
    $\mathsf{Adv}^{\text{SMLWE}}_{\mathcal{A}}(\lambda)=\text{negl}(\lambda)$
    for all PPT adversaries $\mathcal{A}$. Decisional Module-LWE is defined
    similarly to \Cref{def:decision-lwe-assumption}.
\end{definition}

\begin{figure}[ht]
    \centering
    \begin{minipage}{0.45\linewidth}
        \begin{mdframed}[linecolor=black, linewidth=1pt,
                         skipabove=10pt, skipbelow=10pt]
        \colorbox{green!10}{\underline{$\mathsf{MLWEGen}_{\mathcal{R},k,\ell,q,\chi}(\mathbf{s}) \to \mathcal{R}_q^{k \times \ell}
            \times \mathcal{R}_q^{k}$:}}
        \begin{enumerate}
            \item $A \leftarrowS \mathcal{R}_q^{k \times \ell}$;
            \item $\mathbf{e} \leftarrowS \chi$ (over $\mathcal{R}_q^k$);
            \item $\mathbf{b} \gets A\mathbf{s} + \mathbf{e}$;
            \item \textbf{return} $(A,\mathbf{b})$.
        \end{enumerate}
        \end{mdframed}
        \caption{Modular-LWE $(\mathcal{R},k,\ell,q,\chi)$ instance generation} 
        \label{fig:mlwe-gen}
    \end{minipage}
    \hfill
    \begin{minipage}{0.45\linewidth}
        \begin{mdframed}[linecolor=black, linewidth=1pt,
                         skipabove=10pt, skipbelow=10pt]
        \colorbox{green!10}{\underline{$\mathsf{Game}\;\mathsf{SMLWE}_{\mathsf{MLWEGen}}^{\mathcal{A}}(1^{\lambda})$:}}
        \begin{enumerate}
            \item $\mathbf{s} \leftarrowS \mathcal{R}_q^{\ell}$;
            \item $(A,\mathbf{b}) \leftarrowS \mathsf{MLWEGen}(\mathbf{s})$;
            \item $\widehat{\mathbf{s}} \gets \mathcal{A}(A,\mathbf{b})$;
            \item \textbf{return} $\mathbf{s} =_? \widehat{\mathbf{s}}$.
        \end{enumerate}
        \end{mdframed}
        \caption{Search Modular-LWE Game}
        \label{fig:mlwe-game}
    \end{minipage}
\end{figure}

At first glance, it might seem that Module-LWE looks much less efficient than
Ring-LWE: for instance, why should representing $A$ as a matrix
$\mathcal{R}_q^{k \times \ell}$ be any more efficient than using a single
$\mathcal{R}_q$? Notice that to adjust the security of Ring-LWE, one must
regulate $n$ --- the power of cyclotomic polynomial. However, it must be a power
of two, which makes such parameter adjusting highly problematic. For Module-LWE,
one typically fixes $(q,n)$ (which fully determine $\mathcal{R}_q$-arithmetic
efficiency), and then changes $\ell$ to set target security level.

This way, key sizes (which are primarily determined by a storage of $\mathbf{b}$
by Remark~\ref{remark:seeded-a}) turn out to be more or less equal for both
Ring-LWE and Module-LWE:
\begin{itemize}
    \item \textit{Dilithium}~\cite{dilithium}, being a Module-LWE, uses
    parameters $q=2^{23}-2^{13}+1$, $n=256$, and $(k,\ell) \in
    \{(3,2),(4,3),(5,4),(6,5)\}$; If we take, say $(k,\ell)=(4,3)$, we need $kn
    = 1024$ elements of $\mathbb{Z}_q$ to store $\mathbf{b} \in
    \mathcal{R}_q^{k}$.
    \item \textit{NewHope}~\cite{new-hope}, being a Ring-LWE, uses
    $q=2^{13}+2^{12}+1$ and $n=1024$. Notice that although we have a single
    polynomial $b \in \mathcal{R}_q$, we still need $1024$ elements of
    $\mathbb{Z}_q$ for storage.
\end{itemize}

Turns out that implementing analogous PKE scheme to ones in \Cref{def:lwe-pke}
and \Cref{def:rlwe-pke} based on Module-LWE assumptions leads to the Kyber PKE
description! We specify the details in \Cref{section:kyber}.

\section{Lattice-based Cryptographic Protocols}

In \Cref{section:lwe}, we have defined the public-key encryption scheme based on
LWE assumptions (specifically, based on plain-LWE and Ring-LWE). In this
section, we shall see how to implement several familiar cryptographic protocols
using these assumptions.

\subsection{Kyber Public-Key Encryption}\label{section:kyber}

As previously announced, Kyber PKE Scheme~\cite{kyber} is, in the essence, a
Module-LWE-based adaptation of \Cref{def:lwe-pke} and \Cref{def:rlwe-pke}.

\begin{definition}\label{def:kyber} Based on Module-LWE instance with parameters
    $(\mathcal{R},k,\ell,q,\chi)$, define the PKE scheme
    $\Pi_{\text{PKE}}^{\text{Kyber}}$ over message space $\{0,1\}^n$ as follows:
    $\mathsf{KeyGen}$ is specified in \Cref{fig:kyber-keygen}, $\mathsf{Enc}$ in
    \Cref{fig:kyber-enc}, and $\mathsf{Dec}$ in \Cref{fig:kyber-dec}.

    In particular, original paper \cite{kyber} proposes the following
    parameters: $n=256$, $k=\ell=3$, $q=7681$, and $\chi$ is the centered
    binomial distribution $B_{\eta}$ for $\eta=4$.
\end{definition}

\begin{figure}[ht]
\centering
\footnotesize 

\begin{minipage}[t]{0.32\linewidth}
    \begin{mdframed}[linecolor=black, linewidth=1.0pt,
                     skipabove=5pt, skipbelow=5pt,
                     innerleftmargin=5pt, innerrightmargin=5pt]
        \centering
        \colorbox{purple!10}{\underline{$\mathsf{KeyGen}(1^{\lambda}) \to (\mathsf{pk},\mathsf{sk})$:}}
        \begin{enumerate}
            \item $\mathbf{s},\mathbf{e} \leftarrowS \chi^k$ (over $\mathcal{R}_q^k$);
            \item $A \leftarrowS \mathcal{R}_q^{k \times k}$;
            \item $\mathbf{b} \gets A\mathbf{s} + \mathbf{e}$;
            \item \textbf{return} $((A,\mathbf{b}),\mathbf{s})$.
        \end{enumerate}
    \end{mdframed}
    \caption{Kyber Key Generation (note that $\ell=k$)}
    \label{fig:kyber-keygen}
\end{minipage}
\hfill
\begin{minipage}[t]{0.32\linewidth}
    \begin{mdframed}[linecolor=black, linewidth=1.0pt,
                     skipabove=5pt, skipbelow=5pt,
                     innerleftmargin=5pt, innerrightmargin=5pt]
        \centering
        \colorbox{green!10}{\underline{$\mathsf{Enc}(\mathsf{pk}, (m_i)_{i \in (n)} \in \{0,1\}^n) \to c$:}}
        \begin{enumerate}
            \item $m \gets \sum_{i \in (n)}m_iT^i \in \mathcal{R}_q$;
            \item $\Delta \mathbf{r}, \Delta \mathbf{z} \leftarrowS \chi^k$;
            \item $\Delta z' \leftarrowS \chi$;
            \item $\widetilde{\mathbf{a}} \gets A^{\top}\Delta\mathbf{r} + \Delta\mathbf{z}$;
            \item $\widetilde{b} \gets \mathbf{b}^{\top}\Delta \mathbf{r} + \Delta z' + \lceil \frac{q}{2} \rfloor m$;
            \item \textbf{return} $c=(\widetilde{\mathbf{a}}, \widetilde{b})$.
        \end{enumerate}
    \end{mdframed}
    \caption{Kyber Encryption} 
    \label{fig:kyber-enc}
\end{minipage}
\hfill
\begin{minipage}[t]{0.32\linewidth}
    \begin{mdframed}[linecolor=black, linewidth=1.0pt,
                     skipabove=5pt, skipbelow=5pt,
                     innerleftmargin=5pt, innerrightmargin=5pt]
        \centering
        \colorbox{green!10}{\underline{$\mathsf{Dec}(\mathsf{sk}, c) \to m$:}}
        \begin{enumerate}
            \item Parse $(\widetilde{\mathbf{a}}, \widetilde{b}) \gets c$;
            \item $\delta \gets \widetilde{b} - \mathbf{s}^{\top}\widetilde{\mathbf{a}}$;
            \item \textbf{return} $1$ iff $\|\delta\|_{\infty} > \frac{q}{4}$.
        \end{enumerate}
    \end{mdframed}
    \caption{Kyber Decryption}
    \label{fig:kyber-dec}
\end{minipage}
\end{figure}

\textbf{Short-Secret LWE (ss-LWE)*.} Before stating the security of the Kyber
PKE, notice the subtle difference from our previous PKE schemes: here, the
secret value $\mathbf{s}$ is sampled from the same distribution as the error
$\mathbf{e}$. This surely optimizes the storage requirements, but it might seem
that this significantly reduces the security: for starters, the number of
possible values of $\mathbf{s}$ reduces from $q^{nk}$ down to $\eta^{nk}$ (where
each coefficient $c_i$ from output of $\chi$ is bounded by $\eta$:
$\|c_i\|_{\infty} \leq \eta$). However, $\eta^{nk}$ is still an exponentially
large value: even for $\eta=2$, the number of possible values is $2^{768}$ given
Kyber parameters. This motivates us to modify the LWE assumptions to allow
$\mathbf{s}$ to be sampled from smaller distributions. Further, we will show
results for plain LWE, and the following results can be trivially generalized to
Ring and Module LWEs.

\begin{definition}[ss-LWE]
    Let $\underline{\mathsf{Game} \;
    \text{\textsf{ss}-\textsf{LWE}}}_{\mathsf{LWEGen}}^{\mathcal{A}}$ be as
    specified in~\Cref{fig:search-lwe}, but $\mathbf{s}$ is sampled from the
    distribution $\chi$. Similarly define the advantage of adversary $\mathcal{A}$
    as follows:
    \begin{equation*}
        \mathsf{Adv}^{\text{ss-LWE}}_{\mathcal{A},\mathsf{LWEGen}}(\lambda) = \text{Pr}\left[\underline{\mathsf{Game} \;
    \text{\textsf{ss}-\textsf{LWE}}}_{\mathsf{LWEGen}}^{\mathcal{A}}(\lambda) = 1\right].
    \end{equation*}

    We say that \textbf{short-secret LWE assumption} holds if
    $\mathsf{Adv}^{\text{ss-LWE}}_{\mathcal{A},\mathsf{LWEGen}}(\lambda)=\text{negl}(\lambda)$
    for all PPT $\mathcal{A}$.
\end{definition}

We claim that ss-LWE and search LWE assumptions are in fact polynomially
equivalent (although with different LWE instances parameters). We introduce the
following proposition.

\begin{proposition}
    The following two statements are true:
    \begin{enumerate}
        \item[(a)] if \textsf{ss}-\textsf{LWE} assumption holds with parameters
        $(n,m,q,\chi)$, then \textsf{LWE} assumption with the same parameters
        $(n,m,q,\chi)$ holds as well;
        \item[(b)] if \textsf{LWE} assumption holds with parameters
        $(n,m,q,\chi)$, then \textsf{ss}-\textsf{LWE} assumption with parameters
        $(n,m-n,q,\chi)$ holds as well.
    \end{enumerate}
\end{proposition}

\textbf{Proof.} \textit{Part (a).} Assume
$\mathsf{Adv}^{\text{SLWE}}_{\mathcal{A},\mathsf{LWEGen}}(\lambda)$ is
non-negligible for some $\mathcal{A}$. We build $\mathcal{B}^{\mathcal{A}}$ that
uses $\mathcal{A}$ as an oracle and breaks the \textsf{ss}-\textsf{LWE}
assumption in \Cref{fig:ss-lwe-to-s-lwe}. Note that $\mathbf{b} +
A\boldsymbol{\delta}$ is distributed uniformly over $\mathbb{Z}_q^m$, thus
$(A,\widetilde{\mathbf{b}})$ is a valid instance of \textsf{SLWE}, which is
needed for a valid $\mathcal{A}$ execution at step 3. Thus,
\begin{equation*}
    \mathsf{Adv}^{\text{SLWE}}_{\mathcal{A},\mathsf{LWEGen}}(\lambda) = \mathsf{Adv}^{\text{ss-LWE}}_{\mathcal{B}^{\mathcal{A}},\mathsf{LWEGen}}(\lambda).
\end{equation*}

\textit{Part (b).} Now assume
$\mathsf{Adv}^{\text{ss-LWE}}_{\mathcal{A},\mathsf{LWEGen}(n,m-n,q,\chi)}(\lambda)$
is non-negligible for some $\mathcal{A}$. We build $\mathcal{B}^{\mathcal{A}}$
that uses $\mathcal{A}$ as an oracle and breaks the \textsf{SLWE} assumption for
parameters $(n,m,q,\chi)$ in \Cref{fig:s-lwe-to-ss-lwe}. The main trick here is
to rearrange the original equation $\mathbf{b}=A\mathbf{s}+\mathbf{e}$ to get
the \textsf{ss-LWE} instance of smaller size. We do precisely that: following
steps 1--5, we obtain $\widetilde{\mathbf{b}} = \widetilde{A}\mathbf{e}_1 +
\mathbf{e}_2$, which is clearly a \textsf{ss-LWE} instance with parameters
$(n,m-n,q,\chi)$, making oracle call at step 6 valid. Notice that solving this
\textsf{ss-LWE} instance provides $\mathcal{B}$ with an error piece
$\mathbf{e}_1$, making it possible to restore $\mathbf{s}$ by solving a linear
equation $\mathbf{b}_1 = A_1\mathbf{s} + \mathbf{e}_1$. Thus,
\begin{equation*}
    \mathsf{Adv}^{\text{ss-LWE}}_{\mathcal{A},\mathsf{LWEGen}(n,m-n,q,\chi)}(\lambda) = \mathsf{Adv}^{\text{SLWE}}_{\mathcal{B}^{\mathcal{A}},\mathsf{LWEGen}(n,m,q,\chi)}(\lambda). \quad \qedwhite
\end{equation*}

\begin{figure}[ht]
    \centering
    \begin{minipage}{0.3\linewidth}
        \begin{mdframed}[linecolor=black, linewidth=1pt,
                         skipabove=10pt, skipbelow=10pt]
        \colorbox{green!10}{\underline{$\mathcal{B}_{\text{ss-LWE}}^{\mathcal{A}_{\text{SLWE}}}(A,\mathbf{b}) \to \mathbf{s}$:}}
        \begin{enumerate}
            \item $\boldsymbol{\delta} \leftarrowS \mathbb{Z}_q^m$;
            \item $\widetilde{\mathbf{b}} \gets \mathbf{b} + A\boldsymbol{\delta}$;
            \item $\widetilde{\mathbf{s}} \gets \mathcal{A}_{\text{SLWE}}(A,\widetilde{\mathbf{b}})$;
            \item $\widehat{\mathbf{s}} \gets \widetilde{\mathbf{s}} - \boldsymbol{\delta}$
            \item \textbf{return} $\widehat{\mathbf{s}}$.
        \end{enumerate}
        \end{mdframed}
        \caption{Solving ss-LWE given Search-LWE solver} 
        \label{fig:ss-lwe-to-s-lwe}
    \end{minipage}
    \hfill
    \begin{minipage}{0.65\linewidth}
        \begin{mdframed}[linecolor=black, linewidth=1pt,
                         skipabove=10pt, skipbelow=10pt]
        \colorbox{green!10}{\underline{$\mathcal{B}_{\text{SLWE}}^{\mathcal{A}_{\text{ss-LWE}}}(A,\mathbf{b}) \to \mathbf{s}$:}}
        \begin{enumerate}
            \item let $A = \begin{bmatrix}
                A_1 \\ A_2
            \end{bmatrix}$ where $A_1 \in \mathbb{Z}_q^{n \times n}$, $A_2 \in \mathbb{Z}_q^{(m-n) \times n}$;
            \item let $\mathbf{b} = \begin{bmatrix}
                \mathbf{b}_1 \\ \mathbf{b}_2
            \end{bmatrix}$ where $\mathbf{b}_1 \in \mathbb{Z}_q^{n}$, $\mathbf{b}_2 \in \mathbb{Z}_q^{m-n}$;
            \item let $\mathbf{e} = \begin{bmatrix}
                \mathbf{e}_1 \\ \mathbf{e}_2
            \end{bmatrix}$ where $\mathbf{e}_1 \in \mathbb{Z}_q^{n}$, $\mathbf{e}_2 \in \mathbb{Z}_q^{m-n}$;
            \item $\widetilde{A} \gets -A_2A_1^{-1} \in \mathbb{Z}_q^{(m-n) \times n}$;
            \item $\widetilde{\mathbf{b}} \gets \widetilde{A}\mathbf{b}_1 + \mathbf{b}_2 \in \mathbb{Z}_q^{m-n}$;
            \item $\widehat{\mathbf{e}}_1 \gets \mathcal{A}_{\text{ss-LWE}}(\widetilde{A},\widetilde{\mathbf{b}})$;
            \item $\widehat{\mathbf{s}} \gets A_1^{-1}(\mathbf{b}_1 - \mathbf{e}_1)$;
            \item \textbf{return} $\widehat{\mathbf{s}}$.
        \end{enumerate}
        \end{mdframed}
        \caption{Solving Search-LWE given ss-LWE solver}
        \label{fig:s-lwe-to-ss-lwe}
    \end{minipage}
\end{figure}

Now with the ss-LWE assumption in mind, we can state the security of the Kyber
PKE.
\begin{proposition}
    $\Pi_{\text{PKE}}^{\text{Kyber}}$ is correct and IND-CPA secure under
    Short-Secret Decisional Module-LWE assumption.
\end{proposition}

\subsection{Dilithium Digital Signature}\label{section:dilithium}

In this section, we specify the main idea behind the well-known lattice-based
signature scheme \textit{Dilithium}~\cite{dilithium}. Compared to previous
discussion, we omit formal proofs of security and concentrate on the
construction. 

First, recall the definition of the signature scheme.
\begin{definition}[Signature Scheme]
    A \textbf{signature scheme} $\Sigma$ is a triplet of the following effective
    probabilistic algorithms:
    \begin{itemize}
        \item $\mathsf{KeyGen}(1^{\lambda}) \to (\mathsf{pk},\mathsf{sk})$
        outputs the pair of public and secret key, respectively.
        \item $\mathsf{Sign}(\mathsf{sk}, m) \to \sigma$ signs a message 
        $m$ using secret key $\mathsf{sk}$ and thereby produces signature $\sigma$.
        \item $\mathsf{Verify}(\mathsf{pk}, m, \sigma) \to \{0,1\}$ checks
        whether the signature $\sigma$ was produced by signing message $m$ using
        $\mathsf{sk}$ corresponding to $\mathsf{pk}$ and outputs the decision.
    \end{itemize}

    We require \textbf{correctness}:
    \begin{equation*}
        \forall m: \; \text{Pr}\left[ \mathsf{Verify}(\mathsf{pk}, m, \sigma) = 1 \;\middle|\; \begin{aligned}
            &(\mathsf{pk}, \mathsf{sk}) \gets \mathsf{KeyGen}(1^{\lambda}) \\
            &\sigma \gets \mathsf{Sign}(\mathsf{sk}, m)
        \end{aligned} \right] \geq 1 - \text{negl}(\lambda),
    \end{equation*}
    and unforgeability against chosen-message attacks, roughly meaning that
    obtaining $\text{poly}(\lambda)$ pairs $(\sigma_i,m_i)$ signed by some
    key-pair $(\mathsf{pk},\mathsf{sk})$ does now allow $\mathcal{A}$ to produce
    a valid pair $(\widehat{\sigma},\widehat{m})$. See \cite[Section
    13]{BonehShoup2023} for more details.
\end{definition}

\subsubsection{Schnorr Signature}

Further denote by $H$ the hash function from $\{0,1\}^*$ to the appropriate
domain. One of the most widely used digital signatures schemes (especially in
the Bitcoin applications) is the \textit{Schnorr Signature Scheme}. We define it
below.

\begin{definition}\label{def:schnorr}
    Fix a cyclic abelian group $\mathbb{G} = \langle g \rangle$ of order $q$.
    Define the \textbf{Schnorr Signature} scheme $\Sigma_{\text{Schnorr}}$ as a
    triplet of functions specified in \Cref{fig:schnorr-keygen},
    \Cref{fig:schnorr-sign}, and \Cref{fig:schnorr-verify}.
\end{definition}

\begin{figure}[ht]
\centering
\footnotesize 

\begin{minipage}[t]{0.32\linewidth}
    \begin{mdframed}[linecolor=black, linewidth=1.0pt,
                     skipabove=5pt, skipbelow=5pt,
                     innerleftmargin=5pt, innerrightmargin=5pt]
        \centering
        \colorbox{purple!10}{\underline{$\mathsf{KeyGen}(1^{\lambda}) \to (\mathsf{pk},\mathsf{sk})$:}}
        \begin{enumerate}
            \item $\alpha \leftarrowS \mathbb{Z}_q$;
            \item $h \leftarrow g^{\alpha} \in \mathbb{G}$;
            \item \textbf{return} $(h, \alpha)$.
        \end{enumerate}
    \end{mdframed}
    \caption{Schnorr keypair generation}
    \label{fig:schnorr-keygen}
\end{minipage}
\hfill
\begin{minipage}[t]{0.32\linewidth}
    \begin{mdframed}[linecolor=black, linewidth=1.0pt,
                     skipabove=5pt, skipbelow=5pt,
                     innerleftmargin=5pt, innerrightmargin=5pt]
        \centering
        \colorbox{green!10}{\underline{$\mathsf{Sign}(\alpha, m \in \{0,1\}^*) \to \sigma$:}}
        \begin{enumerate}
            \item $r \leftarrowS \mathbb{Z}_q$, $a \gets g^r$;
            \item $e \gets H(m, a)$;
            \item $z \gets r + \alpha e$;
            \item \textbf{return} $\sigma := (e, z)$.
        \end{enumerate}
    \end{mdframed}
    \caption{Schnorr signing} 
    \label{fig:schnorr-sign}
\end{minipage}
\hfill
\begin{minipage}[t]{0.32\linewidth}
    \begin{mdframed}[linecolor=black, linewidth=1.0pt,
                     skipabove=5pt, skipbelow=5pt,
                     innerleftmargin=5pt, innerrightmargin=5pt]
        \centering
        \colorbox{green!10}{\underline{$\mathsf{Verify}(h, m, \sigma) \to \{0,1\}$:}}
        \begin{enumerate}
            \item parse $(e, z) \gets \sigma$;
            \item $a' \gets g^{z}h^{-e}$;
            \item \textbf{return} $1$ iff $e = H(m,a')$.
        \end{enumerate}
    \end{mdframed}
    \caption{Schnorr signature verification}
    \label{fig:schnorr-verify}
\end{minipage}
\end{figure}

Why all of a sudden we decided to bring up the Schnorr signature when
considering LWE-based signature? Turns out that in the essence, the structure of
Dilithium greatly resembles Schnorr signatures! Note that Schnorr signature is a
compiled interactive $\Sigma$-protocol. In the terminology of interactive
protocols, $a$ in a \textit{commitment}, $e$ is a \textit{challenge}, and $z$ is
a \textit{response}. Our hope is that we can construct a similar protocol that
will have different yet resembling expressions for $(a,e,z)$, but in the ``LWE
language''.

\subsubsection{LWE-like $\Sigma$-protocol}\label{section:lwe-sigma}

Fix $(\mathcal{R},k,\ell,q)$ Module-LWE instance. We are going to require a
couple of additional sets to make our scheme work: 
\begin{itemize}
    \item $\mathcal{Q}_{\eta}$ is a set of polynomials $\sum_{j \in (n)}\alpha_jT^j$ in $\mathcal{R}_q$ with $\|\alpha_j\|_{\infty} \leq \eta$.
    \item $\widetilde{\mathcal{Q}}_{\gamma}$ is a set of polynomials $\sum_{j \in
    (n)}\alpha_jT^j$ in $\mathcal{R}_q$ with $\|\alpha_j\|_{\infty} < \gamma$
    for relatively large $\gamma$.
    \item $\mathcal{B}_{\tau}$ is a set polynomials in $\mathcal{Q}_1$, exactly $\tau$
    of whose coefficients are $\pm 1$.
\end{itemize}

Additionally set $\nu := \tau\eta$. For instance, for NIST FIPS 204
standard~\cite{fips-204} one chooses parameters $q=2^{23}-2^{13}+1$, $n=256$,
$(k,\ell)=(8,7)$, $\eta=2$, $\gamma=2^{19}$, $\tau=60$, $\nu = 120$. 

\textbf{First Attempt.} Fix hash function $H: \{0,1\}^* \to \mathcal{B}_{\tau}$.
We present the first attempt of building LWE-based Schnorr-like signature in
\Cref{fig:first-attempt-keygen}, \Cref{fig:first-attempt-sign}, and
\Cref{fig:first-attempt-verify}.

\begin{figure}[ht]
\centering
\footnotesize 

\begin{minipage}[t]{0.32\linewidth}
    \begin{mdframed}[linecolor=black, linewidth=1.0pt,
                     skipabove=5pt, skipbelow=5pt,
                     innerleftmargin=5pt, innerrightmargin=5pt]
        \centering
        \colorbox{purple!10}{\underline{$\mathsf{KeyGen}(1^{\lambda}) \to (\mathsf{pk},\mathsf{sk})$:}}
        \begin{enumerate}
            \item $A \leftarrowS \mathcal{R}_q^{k \times \ell}$;
            \item $\mathbf{s}_1 \leftarrowS \mathcal{Q}_{\eta}^{\ell}, \mathbf{s}_2 \leftarrowS \mathcal{Q}_{\eta}^{k}$;
            \item $\mathbf{t} \gets A\mathbf{s}_1 + \mathbf{s}_2$;
            \item $\mathsf{pk} \gets (A,\mathbf{t}) \in \mathcal{R}_q^{k \times \ell} \times \mathcal{R}_q^{k}$;
            \item $\mathsf{sk} \gets (\mathbf{s}_1,\mathbf{s}_2) \in \mathcal{Q}_{\eta}^{\ell+k}$;
            \item \textbf{return} $(\mathsf{pk},\mathsf{sk})$.
        \end{enumerate}
    \end{mdframed}
    \caption{First-attempt LWE: key generation}
    \label{fig:first-attempt-keygen}
\end{minipage}
\hfill
\begin{minipage}[t]{0.32\linewidth}
    \begin{mdframed}[linecolor=black, linewidth=1.0pt,
                     skipabove=5pt, skipbelow=5pt,
                     innerleftmargin=5pt, innerrightmargin=5pt]
        \centering
        \colorbox{green!10}{\underline{$\mathsf{Sign}(\mathsf{sk}, m \in \{0,1\}^*) \to \sigma$:}}
        \begin{enumerate}
            \item parse $(\mathbf{s}_1,\mathbf{s}_2) \gets \mathsf{sk}$;
            \item $\mathbf{r} \leftarrowS \widetilde{\mathcal{Q}}_{\gamma}^{\ell}$;
            \item $\mathbf{a} \gets A\mathbf{r}$; \hspace{3px}{\textcolor{gray!70!black}{\footnotesize\textit{// commitment}}}
            \item $e \gets H(m,\mathbf{a})$; \hspace{3px}{\textcolor{gray!70!black}{\footnotesize\textit{// challenge}}}
            \item $\mathbf{z} \gets \mathbf{r} + e\mathbf{s}_1$; \hspace{3px}{\textcolor{gray!70!black}{\footnotesize\textit{// response}}}
            \item $\sigma \gets (e,\mathbf{z}) \in \mathcal{B}_{\tau} \times \mathcal{R}_q^{\ell}$;
            \item \textbf{return} $\sigma$.
        \end{enumerate}
    \end{mdframed}
    \caption{First-attempt LWE: signing} 
    \label{fig:first-attempt-sign}
\end{minipage}
\hfill
\begin{minipage}[t]{0.32\linewidth}
    \begin{mdframed}[linecolor=black, linewidth=1.0pt,
                     skipabove=5pt, skipbelow=5pt,
                     innerleftmargin=5pt, innerrightmargin=5pt]
        \centering
        \colorbox{green!10}{\underline{$\mathsf{Verify}(\mathsf{pk}, m, \sigma) \to \{0,1\}$:}}
        \begin{enumerate}
            \item parse $(e, \mathbf{z}) \gets \sigma$;
            \item parse $(A,\mathbf{t}) \gets \mathsf{pk}$;
            \item compute $\mathbf{a}'$;
            \item \textbf{return} $1$ iff $e=H(m,\mathbf{a}')$.
        \end{enumerate}
    \end{mdframed}
    \caption{First-attempt LWE: signature verification}
    \label{fig:first-attempt-verify}
\end{minipage}
\end{figure}

Notice that the protocol greatly resembles the Schnorr signature scheme
previously presented in \Cref{def:schnorr}. However, the biggest issue so far is
step 3 in \Cref{fig:first-attempt-verify}: how can the verifier recompute the
commitment $\mathbf{a}'$ specifically?

Note that $\mathbf{z} = \mathbf{r} + e\mathbf{s}_1$, thus $A\mathbf{z} =
A\mathbf{r} + eA\mathbf{s}_1 = \mathbf{a} + e(\mathbf{t}-\mathbf{s}_2)$ and so
we obtain $A\mathbf{z} - e\mathbf{t} = \mathbf{a} - e\mathbf{s}_2$.
Unfortunately, this way we only obtain $\mathbf{a}-e\mathbf{s}_2$ where
$\mathbf{s}_2$ is unknown to the verifier. However, here where our distributions
$\mathcal{Q}_{\eta}$ and $\mathcal{B}_{\tau}$ might come into play: notice that
$e\mathbf{s}_2$ is a vector of polynomials in $\mathcal{R}_q^k$ with small
coefficients (since $\|s_{2,i}\|_{\infty} \leq \eta$ for all $i \in [k]$ and
$\|e\|_{\infty} \leq 1$ by construction).

Thus, we might modify our idea in the following way: the signer will use the
function called $\mathsf{HighBits}$, which will divide and round each
coefficient of $\mathbf{w}$ which hopefully eventually removes the noise
$e\mathbf{s}_2$ (we give concrete specifics later on). The signer then composes
the challenge $e = H(m,\mathbf{a}_H)$ where $\mathbf{a}_H =
\mathsf{HighBits}(\mathbf{a})$.

If $\mathsf{LowBits}$ of $\mathbf{a}-e\mathbf{s}_2$ is sufficiently small then,
as $e\mathbf{s}_2$ is small, we have
$\mathsf{HighBits}(\mathbf{a}-e\mathbf{s}_2)=\mathsf{HighBits}(\mathbf{a})$, and
so the verifier can compute $\mathbf{a}'$. This way, we modify the signing
procedure by making signer recompute $(e,\mathbf{z})$ while
$\mathsf{LowBits}(\mathbf{a}-e\mathbf{s}_2)$ is not sufficiently small.

\textbf{Second Attempt.} We encapsulate observations above in
\Cref{fig:second-attempt-sign} and \Cref{fig:second-attempt-verify}.

\begin{figure}[ht]
\centering
\begin{minipage}[t]{0.45\linewidth}
    \begin{mdframed}[linecolor=black, linewidth=1.0pt,
                     skipabove=5pt, skipbelow=5pt,
                     innerleftmargin=5pt, innerrightmargin=5pt]
        \centering
        \colorbox{green!10}{\underline{$\mathsf{Sign}(\mathsf{sk}, m \in \{0,1\}^*) \to \sigma$:}}
        \begin{enumerate}
            \item parse $(\mathbf{s}_1,\mathbf{s}_2) \gets \mathsf{sk}$;
            \item \textcolor{green!50!black}{\textbf{while} $\mathsf{LowBits}(\mathbf{a}-e\mathbf{s}_2)$ is large}:
            \item \hspace{20px}$\mathbf{r} \leftarrowS \widetilde{\mathcal{Q}}_{\gamma}^{\ell}$;
            \item \hspace{20px}$\mathbf{a} \gets A\mathbf{r}$; 
            \item \hspace{20px}\textcolor{green!50!black}{$\mathbf{a}_H \gets \mathsf{HighBits}(\mathbf{a})$};
            \item \hspace{20px}$e \gets H(m, \textcolor{green!50!black}{\mathbf{a}_H})$; 
            \item \hspace{20px}$\mathbf{z} \gets \mathbf{r} + e\mathbf{s}_1$;
            \item $\sigma \gets (e,\mathbf{z})$;
            \item \textbf{return} $\sigma$.
        \end{enumerate}
    \end{mdframed}
    \caption{Second-attempt LWE: signing} 
    \label{fig:second-attempt-sign}
\end{minipage}
\hfill
\begin{minipage}[t]{0.4\linewidth}
    \begin{mdframed}[linecolor=black, linewidth=1.0pt,
                     skipabove=5pt, skipbelow=5pt,
                     innerleftmargin=5pt, innerrightmargin=5pt]
        \centering
        \colorbox{green!10}{\underline{$\mathsf{Verify}(\mathsf{pk}, m, \sigma) \to \{0,1\}$:}}
        \begin{enumerate}
            \item parse $(e, \mathbf{z}) \gets \sigma$;
            \item parse $(A,\mathbf{t}) \gets \mathsf{pk}$;
            \item \textcolor{green!50!black}{$\mathbf{a}'_H \gets \mathsf{HighBits}(A\mathbf{z} - e\mathbf{t})$};
            \item \textbf{return} $1$ iff $e=H(m,\textcolor{green!50!black}{\mathbf{a}'_H})$.
        \end{enumerate}
    \end{mdframed}
    \caption{Second-attempt LWE: signature verification}
    \label{fig:second-attempt-verify}
\end{minipage}
\end{figure}

However,
this attempt also contains numerous problems we have to mitigate:
\begin{itemize}
    \item \textbf{Issue \#1.} Public key $(A,\mathbf{t}) \in \mathcal{R}_q^{k
    \times \ell} \times \mathcal{R}_q^k$ is very large: we are going to derive
    $A$ from the public seed $\rho$ and use only ``high order bits'' of
    coefficients in $\mathbf{t}$.
    \item \textbf{Issue \#2.} Secret key $(\mathbf{s}_1,\mathbf{s}_2) \in
    \mathcal{Q}_{\eta}^{\ell+k}$ is similarly too large: we are going to
    generate these values from the \textit{secret seed $\rho'$}.
    \item \textbf{Issue \#3.} Costly hashing: during signing, on each step we
    have to perform $H(m,\mathbf{a}_H)$, which is expensive if $m$ is large.
    Instead, we are going to set $\mu \gets H(m)$ and on each step compute
    $H(\mu,\mathbf{a}_H)$.
    \item \textbf{Issue \#4.} Scheme is non-deterministic (which might be
    undesirable for some applications) and requires too many random samplings
    $\mathbf{r} \leftarrowS \widetilde{\mathcal{Q}}_{\gamma}^{\ell}$. Instead,
    we will generate $\mathbf{r}$ from the seed $\rho''$ and counter $\kappa$,
    where $\rho''$ in turn will be obtained by hashing a secret seed $K$ and
    $\mu$.
    \item \textbf{Issue \#5.} This protocol is not fully zero-knowledge: by
    careful inspection of coefficients of $\mathbf{z}$, one might sometimes get
    some information about these coefficients (e.g., if some coefficient is
    $\gamma+\nu$, then the corresponding coefficient of $e\mathbf{s}_1$ is
    $\nu$). The solution is that signer, while producing the signature, will
    ensure that $\|\mathbf{z}\|_{\infty} < \gamma-\nu$.
\end{itemize}

\subsubsection{Dilithium Protocol Simplified Specification}

\textbf{HighBits and LowBits.} Now, let us define these functions 
more carefully.
\begin{definition}
    Fix even $\alpha \in \mathbb{N}$ such that $q-1=m\alpha$. Functions
    $\mathsf{HighBits}_{\alpha}(x)$ and $\mathsf{LowBits}_{\alpha}(x)$ will be
    informally defined as quotient and remainder of division $x$ by $\alpha$,
    respectively. More formally, let $x = x_H\alpha+x_L$ where $|x_L| \leq
    \alpha/2$ and $0 \leq x_H \leq m$. Then, we set
    $\mathsf{HighBits}_{\alpha}(x) := x_H$ and $\mathsf{LowBits}_{\alpha}(x) :=
    x_L$.
\end{definition}

\textbf{Additional parameters.} We slighly modify parameters specified earlier
in \Cref{section:lwe-sigma}: instead of $\gamma$, we use $\gamma_1$, and we
additionally introduce parameters $\gamma_2$ (which in NIST FIPS 204, for
instance, is $(q-1)/32 \approx 2^{18}$) and $\alpha := 2\gamma_2$. Additionally,
denote by $H(\bullet;d): \{0,1\}^* \to \{0,1\}^d$ the hash function returning
$d$-bit binary string as an output, $H_{\mathcal{B}}(\bullet): \{0,1\}^{2\lambda} \to
\mathcal{B}_{\tau}$ the hash function returning polynomial in
$\mathcal{B}_{\tau}$, and by $\gets_{\rho}$ sampling with seed $\rho$.

Finally, we are ready to present almost complete specification of the Dilithium
protocol.

\begin{definition}[Dilithium Signature Scheme]
    We define the \textbf{Dilithium} signature scheme
    $\Sigma_{\text{Dilithium}}$ as a triplet of functions specified in
    \Cref{fig:delithium-keygen}, \Cref{fig:delithium-sign}, and
    \Cref{fig:delithium-verify}.
\end{definition}

Let us again check the \textit{correctness} of the obtained scheme. Notice that
\begin{equation*}
    A\mathbf{z} - e\mathbf{t} = A(\mathbf{r}+e\mathbf{s}_1) - e(A\mathbf{s}_1+\mathbf{s}_2) = A\mathbf{r} - e\mathbf{s}_2 = \mathbf{a} - e\mathbf{s}_2.
\end{equation*}

Since $\|\mathsf{LowBits}_{2\gamma_2}(\mathbf{a}-e\mathbf{s}_2)\|_{\infty} \leq
\gamma_2 - \nu$ and $\|e\mathbf{s}_2\|_{\infty} \leq \nu$, we have
\begin{equation*}
    \mathbf{a}_H' = \mathsf{HighBits}_{2\gamma_2}(A\mathbf{z}-e\mathbf{t}) = \mathsf{HighBits}_{2\gamma_2}(\mathbf{a}-e\mathbf{s}_2) = \mathsf{HighBits}_{2\gamma_2}(\mathbf{a}) = \mathbf{a}_H.
\end{equation*}

Thus, valid signatures are accepted.

\begin{figure}[H]
\centering
\begin{varwidth}{0.75\linewidth}
\begin{mdframed}[linecolor=black, linewidth=1.0pt,
                    skipabove=5pt, skipbelow=5pt,
                    innerleftmargin=5pt, innerrightmargin=5pt]
\centering
\colorbox{purple!10}{\underline{$\mathsf{KeyGen}(1^{\lambda}) \to (\mathsf{pk},\mathsf{sk})$:}}
\begin{enumerate}
    \item $(\rho,\rho',K) \leftarrowS \{0,1\}^{n} \times \{0,1\}^{2n} \times \{0,1\}^{n}$; \hspace{5px}{\textcolor{gray!70!black}{\footnotesize\texttt{// Generate seeds}}}
    \item $A \leftarrowS_{\rho} \mathcal{R}_q^{k \times \ell}$;\hspace{5px}{\textcolor{gray!70!black}{\footnotesize\texttt{// Sample matrix from seed $\rho$}}}
    \item $(\mathbf{s}_1,\mathbf{s}_2) \leftarrowS_{\rho'} \mathcal{Q}_{\eta}^{\ell}\times\mathcal{Q}_{\eta}^k$;\hspace{5px}{\textcolor{gray!70!black}{\footnotesize\texttt{// Sample secrets from seed $\rho'$}}}
    \item $\mathbf{t} \gets A\mathbf{s}_1+\mathbf{s}_2$;
    \item $\mathsf{tr} \gets H(\rho,\mathbf{t};2\lambda)$;
    \item $\mathsf{pk} \gets (\rho,\mathbf{t})$;\hspace{5px}{\textcolor{gray!70!black}{\footnotesize\texttt{// $\rho$ can be used for restoring $A$}}}
    \item $\mathsf{sk} \gets (\rho',K,\textcolor{gray!70!black}{\mathsf{tr}})$;\hspace{5px}{\textcolor{gray!70!black}{\footnotesize\texttt{// $\rho'$ can be used for restoring $\mathbf{s}_1$ and $\mathbf{s}_2$}}}
    \item \textbf{return} $(\mathsf{pk},\mathsf{sk})$.
\end{enumerate}
\end{mdframed}
\end{varwidth}
\caption{Dilithium key generation}
\label{fig:delithium-keygen}
\end{figure}
\begin{figure}[H]
\centering
\begin{varwidth}{0.85\linewidth}
\begin{mdframed}[linecolor=black, linewidth=1.0pt,
                     skipabove=5pt, skipbelow=5pt,
                     innerleftmargin=5pt, innerrightmargin=5pt]
    \centering
    \colorbox{green!10}{\underline{$\mathsf{Sign}(\mathsf{sk},m) \to \sigma$:}}
    \begin{enumerate}
        \item $A \leftarrowS_{\rho} \mathcal{R}_q^{k \times \ell}$;\hspace{5px}{\textcolor{gray!70!black}{\footnotesize\texttt{// resample same $A$ from seed $\rho$}}}
        \item $\mu \gets H(\mathsf{tr},m; 2n)$;
        \item $\mathsf{rnd} \gets \mathbf{0}^n$ \textbf{if} deterministic mode, $\mathsf{rnd} \leftarrowS \{0,1\}^n$ otherwise;
        \item $\rho'' \gets H(K,\mathsf{rnd},\mu; 2n)$; 
        \item $\kappa \gets 0$; $\mathbf{z} \gets \,\perp$; \hspace{5px}{\textcolor{gray!70!black}{\footnotesize\texttt{// initialize counter}}}
        \item \textbf{while} $\mathbf{z} = \,\perp$:
        \item \hspace{20px} $\mathbf{r} \leftarrowS_{(\rho'',\kappa)}\widetilde{\mathcal{Q}}_{\gamma_1}^{\ell}$;
        \item \hspace{20px} $\mathbf{a} \gets A\mathbf{r}$; $\mathbf{a}_H \gets \mathsf{HighBits}_{2\gamma_2}(\mathbf{a})$;\hspace{5px}{\textcolor{gray!70!black}{\footnotesize\texttt{// compute high-bits of commitment}}}
        \item \hspace{20px} $\widetilde{e} \gets H(\mu,\mathbf{a}_H; 2\lambda)$;
        \item \hspace{20px} $e \gets H_{\mathcal{B}}(\widetilde{e}) \in \mathcal{B}_{\tau}$;\hspace{5px}{\textcolor{gray!70!black}{\footnotesize\texttt{// compute challenge}}}
        \item \hspace{20px} $\mathbf{z} \gets \mathbf{r} + e\mathbf{s}_1$;\hspace{5px}{\textcolor{gray!70!black}{\footnotesize\texttt{// compute response}}}
        \item \hspace{20px} $\bm{\delta} \gets \mathsf{LowBits}_{2\gamma_2}(\mathbf{a} - e\mathbf{s}_2)$;\hspace{5px}{\textcolor{gray!70!black}{\footnotesize\texttt{// compute ``noise'' level}}}
        \item \hspace{20px} \textbf{if} $\|\mathbf{z}\|_{\infty} \geq \gamma_1-\nu$ or $\|\bm{\delta}\|_{\infty} \geq \gamma_2-\nu$, then set $\mathbf{z} \gets \,\perp$;
        \item \hspace{20px} $\kappa \gets \kappa+\ell$;\hspace{5px}{\textcolor{gray!70!black}{\footnotesize\texttt{// update counter}}}
        \item \textbf{return} $\sigma=(\widetilde{e}, \mathbf{z})$.
    \end{enumerate}
\end{mdframed}
\end{varwidth}
\caption{Dilithium message signing}
\label{fig:delithium-sign}
\end{figure}
\begin{figure}[H]
\centering
\begin{varwidth}{0.85\linewidth}
\begin{mdframed}[linecolor=black, linewidth=1.0pt,
                     skipabove=5pt, skipbelow=5pt,
                     innerleftmargin=5pt, innerrightmargin=5pt]
    \centering
    \colorbox{green!10}{\underline{$\mathsf{Verify}(\mathsf{pk},m,\sigma) \to \{0,1\}$:}}
    \begin{enumerate}
        \item parse $(\rho,\mathbf{t}) \gets \mathsf{pk}$;
        \item parse $(\widetilde{e},\mathbf{z}) \gets \sigma$;
        \item \textbf{if} $\|\mathbf{z}\|_{\infty} \geq \gamma_1-\nu$ then \textbf{return} 0;
        \item $A \leftarrowS_{\rho} \mathcal{R}_q^{k \times \ell}$;
        \item $\mathsf{tr} \gets H(\rho,\mathbf{t};2n)$; $\mu \gets H(\mathsf{tr},m;2n)$;
        \item $e \gets H_{\mathcal{B}}(\widetilde{e}) \in \mathcal{B}_{\tau}$;
        \item $\mathbf{a}_H' \gets \mathsf{HighBits}_{2\gamma_2}(A\mathbf{z}-e\mathbf{t})$;
        \item \textbf{return} $1$ \textbf{if} $\widetilde{e} = H(\mu,\mathbf{a}_H';2\lambda)$ and $0$ otherwise.
    \end{enumerate}
\end{mdframed}
\end{varwidth}
\caption{Dilithium signature verification}
\label{fig:delithium-verify}
\end{figure}

Finally, we state the security of $\Sigma_{\text{Dilithium}}$ without a proof.
Curious reader is of course encouraged to read the original paper
\cite{dilithium} for details.
\begin{proposition}
    $\Sigma_{\text{Dilithium}}$ is \textbf{zero-knowledge} and
    \textbf{unforgeable against chosen-message} attack assuming (1) Decisional
    Module-LWE, (2) Module-SIS (Short-Integer Solution)\footnote{We have not
    defined it rigorously, but informally we assume finding a \textit{short}
    solution to $[A|E_{k \times k}]\mathbf{s}=0$ for $A \in \mathcal{R}_q^{k
    \times \ell}$ and identity matrix $E_{k \times k}$ is hard.}, and (3)
    (Quantum) Random Oracle Model (QROM).
\end{proposition}

\subsection*{Acknowledgements}

Some of the arguments and ideas used in these notes were particularly inspired
by awesome courses ``Cryptography 101 with Alfred Menezes''\footnote{See
\url{https://cryptography101.ca/}.} (\Cref{section:dilithium} is primarily based
on ``Kyber and Dilithium'' lectures) and lecture materials by François Charles
on lattices during ICMU Summer School on Number Theory.

\subsection*{References}
\footnotesize
\renewcommand{\refname}{\vspace{-1em}}
\bibliographystyle{alphaurl}
\bibliography{refs-lattices}

\end{document}